\begin{textsample}{POS Dim 1 – 2014\_c6 – Score -7.00 – 2014\_c6/enade\_2014\_letras\_por\_6f9b30ff67.txt}  \label{ex:f1_pos_424}
CAPÍTULO 3 ANÁLISE TÉCNICA DA PROVA Este capítulo tem por objetivo apresentar o desempenho dos estudantes concluintes de Letras - Português ( Bacharelado e Licenciatura ) no ENADE/2014.
Para isso, foram calculadas as estatísticas básicas da prova em seu todo ( seção 3.1.1 ), bem como as estatísticas dos componentes relacionadas ao Componente de Formação Geral ( seção 3.1.2 ) e ao de Conhecimento Específico da Área ( seção 3.1.3 ).
Dadas as suas características, foram analisadas, em separado, as questões objetivas ( seção 3.2 ) e as questões discursivas ( seção 3.3 ) para concluintes de Bacharelado e de Licenciatura.
Para as questões objetivas, foram disponibilizados os índices de facilidade e o de discriminação bisserial, também em separado, para os Componentes de Formação Geral ( seção 3.2.1 ) e de Conhecimento Específico ( seção 3.2.2 ), sendo a análise da segunda seção separada por concluintes de Bacharelado e da Licenciatura.
De cada componente, uma das questões foi escolhida para exemplificar a análise gráfica, relacionando as alternativas escolhidas pelos estudantes ( inclusive o gabarito ) com o número de acertos no componente aferido, como já dito anteriormente, agregados por Bacharelado e Licenciatura.
O Anexo I apresenta a íntegra da análise gráfica para todas as questões objetivas agregadas e também separadas por Bacharelado e Licenciatura, a partir da questão 27.
Para cada uma das questões discursivas, os conteúdos dos tipos mais comuns de respostas dos estudantes são apresentados e comparados com o padrão de respostas esperado ( ver Anexo VIII com o padrão de respostas ).
Tomando como base as duas questões discursivas do Componente de Formação Geral, a seção 3.3.1.6 apresenta comentários sobre a correção das respostas com respeito à Língua Portuguesa.
Na Tabela 3.1 são apresentados o tamanho da população inscrita e de presentes, e as seguintes estatísticas das notas19 : média do desempenho na prova, erro padrão da média, desvio padrão, nota mínima, mediana e nota máxima.
Tais estatísticas contemplam o total de estudantes concluintes da Área de Letras - Português ( Bacharelado e Licenciatura ) inscritos e presentes à prova do ENADE/2014, tendo em vista as agregações por Grandes Regiões e o país como um todo.
As estatísticas das demais tabelas ímpares serão baseadas na mesma população da Tabela 3.1 e, portanto, do mesmo tamanho.
As tabelas pares são desagregadas por Categoria Administrativa e Organização Acadêmica da IES.
As estatísticas dessas tabelas pares serão baseadas na mesma população da Tabela 3.2. 19 Uma definição dessas estatísticas pode ser encontrada no Glossário.
Em relação aos gráficos de distribuição de notas, o intervalo considerado foi de 10 unidades, aberto à esquerda e fechado à direita, com exceção do primeiro intervalo, [ 0 ; 10 ], fechado em ambos os extremos.
Para os gráficos de distribuição das notas das questões discursivas, foram consideradas mais duas categorias : questão em branco20 e nota zero.
Todos os gráficos de distribuição de notas permitem a comparação dos resultados por habilitação : Bacharelado e Licenciatura. 3.1 ESTATÍSTICAS BÁSICAS DA PROVA Esta seção apresenta estatísticas selecionadas e histogramas da nota geral ( 3.1.1 ) e de cada componente : Formação Geral ( 3.1.2 ) e Conhecimento Específico ( 3.1.3 ).
São também apresentadas as estatísticas selecionadas de subpopulações caracterizadas por Grande Região, Categoria Administrativa e Organização Acadêmica. 3.1.1 Estatísticas Básicas Gerais A Tabela 3.1 apresenta as estatísticas básicas da prova, por grande Região, dos estudantes concluintes de Letras - Português ( Bacharelado e Licenciatura ).
A população total de inscritos foi de 18.264.
Destes, 13.642 estiveram presentes, sendo 25,3 \% o índice de não comparecimento.
A Região de maior abstenção foi a Norte ( 42,1 \% ), e a de menor abstenção foi a Nordeste ( 19,1 \% ).
A média das notas da prova como um todo ( nas seções seguintes serão analisados os Componentes de Formação Geral e de Conhecimento Específico ) foi 42,2, sendo que os estudantes da região Norte obtiveram a média mais baixa ( 39,0 ), e os da região Sudeste obtiveram a média mais alta ( 45,5 ).
As demais médias foram : 41,6 nas regiões Nordeste e Centro-Oeste ; e 41,4 na região Sul.
O desvio padrão para o Brasil como um todo foi 13,4, sendo o maior desvio padrão encontrado na região Sudeste ( 14,6 ) e o menor, nas regiões Norte e Nordeste ( 12,5 em ambas ), indicando uma dispersão um pouco menor das notas destas últimas regiões.
A região que obteve a maior nota máxima foi a Nordeste ( 91,1 ), ao passo que a região que atingiu a menor nota máxima foi a Norte ( 80,6 ).
A mediana do Brasil como um todo foi 41,7, sendo a maior mediana obtida na região Sudeste ( 45,5 ), e a menor obtida na Norte ( 38,7 ).
A nota mínima foi zero em todas as regiões, sem exceção.
Considerando -se as notas segundo Grande Região, observa -se que existem diferenças estatisticamente significativas ao nível de 95 \% entre a maior média, obtida na região Sudeste ( 45,5 ), e as médias de todas as demais regiões.
O comportamento das notas dos estudantes de todo o Brasil pode ser observado no Gráfico 3.1 que apresenta um histograma com a distribuição das mesmas em colunas diferentes para estudantes de Bacharelado e de Licenciatura.
As distribuições são unimodais.
O intervalo modal dos concluintes de Bacharelado é o ( 50 ; 60 ], e dos concluintes de Licenciatura é o ( 30 ; 40 ].
A Tabela 3.2 apresenta as estatísticas básicas da prova, desagregadas por Categoria Administrativa e por Organização Acadêmica.
Da população total de inscritos, 5.872 são de IES Privadas e 12.392 de IES Públicas.
Em relação à Organização Acadêmica, a maior participação foi obtida por estudantes de Universidades ( 15.445 ), seguido de Centros Universitários ( 1.490 ) e Faculdades ( 1.329 ).
A Categoria Administrativa de maior abstenção foi a Pública ( 26,3 \% ), e entre as Organizações Acadêmicas foi a dos Centros Universitários ( também 26,3 \% ), ambos acima da média nacional de 25,3 \%.
A média das notas da prova como um todo foi 42,2.
Em relação à Categoria Administrativa, os estudantes das IES Públicas obtiveram média mais alta ( 42,7 ), e os das IES Privadas obtiveram média mais baixa ( 41,0 ), que a média nacional.
Observa -se que existe diferença estatisticamente significativa entre as médias das notas das IES Públicas e Privadas.
A diferença entre as médias das regiões Sudeste e Norte ( 6,5 ), a maior e a menor média, é superior à diferença entre IES Públicas e Privadas ( 1,7 ), caracterizando uma maior diversidade regional do que administrativa.
No tocante à Organização Acadêmica, apenas as Universidades obtiveram média mais alta ( 42,5 ) que a nacional.
As demais médias foram : 40,9 nos Centros Universitários e 40,0 nas Faculdades.
O desvio padrão para as IES Públicas ( 13,7 ) e para as Universidades ( 13,6 ) foi superior ao do Brasil como um todo ( 13,4 ), indicando uma dispersão um pouco maior das notas nesta Categoria Administrativa e nesta Organização Acadêmica.
Constata -se que existe diferença estatisticamente significativa ao nível de 95 \% nas médias das notas dos estudantes provenientes de Universidades em relação aos de Centros Universitários e Faculdades. 3.1.2 Estatísticas Básicas no Componente de Formação Geral A Tabela 3.3 apresenta as estatísticas básicas em relação ao componente da prova que avalia a Formação Geral dos estudantes concluintes.
Os estudantes de todo Brasil obtiveram desempenho médio de 52,9.
Quanto à variabilidade, o desvio padrão das notas dos estudantes do Brasil, como um todo, foi 17,8.
A maior média foi obtida na região Sudeste ( 56,0 ), e a menor, na região Norte ( 49,1 ).
As demais médias foram : 52,5 na região Nordeste ; 52,8 na região Sul ; e 52,2 na região Centro-Oeste.
Já o maior desvio padrão foi obtido na região Sudeste ( 18,7 ), e o menor, na região Nordeste ( 17,0 ).
Os demais desvios padrões foram : 17,7 na região Norte, 17,3 na região Sul e 18,4 na região Centro-Oeste.
A maior nota no Componente de Formação Geral da prova do ENADE/2014 foi 99,0, obtida por, pelo menos, um estudante na região Sudeste.
A menor nota máxima foi obtida na região Norte ( 96,4 ).
A mediana do Brasil, como um todo, foi 53,8, sendo a menor mediana encontrada na região Norte ( 49,3 ) ; e a maior encontrada na região Sudeste ( 57,5 ).
A nota mínima nesta parte foi zero em todas as regiões, sem exceção.
Considerando -se as notas segundo Grande Região, observa -se que existe diferença estatisticamente significativa entre a maior média das notas do Componente de Formação Geral, obtida na região Sudeste ( 56,0 ) e as médias de todas as demais regiões.
O Gráfico 3.2 propicia a avaliação do desempenho dos estudantes no Componente de Formação Geral a partir do histograma da distribuição das notas correspondentes ao Bacharelado e à Licenciatura.
As distribuições são unimodais, com a modas, respectivamente, no intervalos ( 60 ; 70 ] e ( 50 ; 60 ], enquanto na prova como um todo ( Gráfico 3.1 ), a moda foi alcançada no intervalo imediatamente inferior para o Bacharelado, ou seja, ( 50 ; 60 ], e no intervalo imediatamente superior, ( 40 ; 50 ], para a Licenciatura.
Nota -se, ainda, que no Gráfico 3.2 as notas apresentam uma maior dispersão do que no Gráfico 3.1 ( distribuição das notas da prova ), confirmado pela comparação dos desvios padrões : 13,4 para a nota da prova como um todo e 17,8 para o Componente de Formação Geral.
Na Tabela 3.4 são apresentadas as informações referentes ao desempenho dos concluintes do Componente de Formação Geral, em diferentes agregações : Categoria Administrativa e Organização Acadêmica.
Observa -se que não existe diferença estatisticamente significativa entre as médias dos tipos de Categoria Administrativa.
A maior média foi obtida por estudantes de IES Públicas ( 53,0 ), sem uma diferença estatisticamente significativa da obtida por estudantes de IES Privadas ( 52,8 ).
Considerando -se o tipo de Organização Acadêmica, nota -se que não há diferença estatisticamente significativa entre as médias dos estudantes de Centros Universitários ( 52,4 ) e de Faculdades ( 50,6 ) ; assim como entre os Centros Universitários e as Universidades ( 53,2 ).
Já entre as Faculdades e as Universidades, existe diferença estatisticamente significativa.
Universitários e as Universidades ( 53,2 ).
Já entre as Faculdades e as Universidades, existe diferença estatisticamente significativa. 3.1.3 Estatísticas Básicas do Componente de Conhecimento Específico A Tabela 3.5 apresenta as estatísticas básicas referentes ao Componente de Conhecimento Específico da área de Letras - Português ( Bacharelado e Licenciatura ).
A média do desempenho dos estudantes do Brasil, como um todo, foi 38,6.
A maior média foi obtida na região Sudeste ( 42,0 ), e a menor, na região Norte ( 35,6 ).
As demais médias foram : 37,9 na região Nordeste ; 37,6 na região Sul ; e 38,0 na região Centro-Oeste.
Quanto à variabilidade das notas, o desvio padrão do Brasil, como um todo, foi 14,3, sendo o maior desvio padrão observado na região Sudeste ( 15,5 ), e o menor, na região Norte ( 13,1 ).
Os demais desvios foram : 13,5 na região Nordeste, 14,1 na região Sul e 14,2 na região Centro-Oeste.
A mediana das notas dos estudantes de todo o Brasil foi 37,8.
A maior mediana ocorreu na região Sudeste ( 41,8 ), e a menor, na região Norte ( 34,9 ).
As demais medianas foram : 37,4 na região Nordeste ; 36,3 na região Sul ; e 36,8 na região Centro-Oeste.
A nota máxima do Brasil, como um todo, foi 91,2, sendo obtida por, pelo menos, um estudante na região Nordeste.
As demais notas máximas foram : 81,0 na região Norte ; 90,7 na região Sudeste ; e 83,5 na região Sul ; e 88,0 na região Centro-Oeste.
A nota mínima foi zero em todas as regiões, sem exceção.
Observa -se que, novamente, existe diferença estatisticamente significativa da maior média das notas do Componente de Conhecimento Específico, a da região Sudeste ( 42,0 ), em relação às médias das demais regiões.
O intervalo de confiança para a média da região Centro-Oeste ( 2,0 ) é maior do que de todas as outras.
Assim como os Gráficos 3.1 e 3.2, o Gráfico 3.3, apresentado a seguir, permite uma avaliação do desempenho de concluintes dos cursos de Bacharelado e de Licenciatura, desta vez em relação ao Componente de Conhecimento Específico, com um histograma da distribuição das notas correspondentes.
Dentre as três distribuições apresentadas, esta é a mais concentrada nas notas baixas.
Estas também são distribuições unimodais, e o grupo modal é o ( 40 ; 50 ] para o Bacharelado, e o ( 30 ; 40 ], para a Licenciatura.
A Tabela 3.6 apresenta uma comparação dos resultados em relação à Categoria Administrativa e à Organização Acadêmica, agora levando em conta o desempenho de estudantes do Componente de Conhecimento Específico da prova.
No que se refere à Organização Acadêmica, a maior média foi das Universidades ( 38,9 ), vindo a seguir a dos Centros Universitários ( 37,0 ) e, depois, a das Faculdades ( 36,5 ).
O maior desvio padrão, e acima do valor para o Brasil, como um todo, foi o das Universidades ( 14,4 ).
As Universidades também obtiveram as maiores notas máxima ( 91,2 ) e mediana ( 38,3 ), seguidas pelos os Centros Universitários ( 85,2 e 35,8 respectivamente ) e pelas Faculdades ( 80,9 e 34,9 respectivamente ).
A nota mínima foi zero para quase todas as Organizações Acadêmicas, exceto para as Faculdades ( 3,8 ).
Quanto à Categoria Administrativa, observa -se um comportamento semelhante àquele da prova como um todo, ou seja, existe diferença estatisticamente significativa entre as médias das IES Públicas ( 39,3 ) e IES Privadas ( 37,1 ).
Neste caso também, a maior média foi obtida por estudantes de IES Públicas de ensino.
Observa -se que existem diferenças estatisticamente significativas ao nível de 95 \% no Componente de Conhecimento Específico entre as notas das Universidades em relação aos Centros Universitários e Faculdades. 3.2 ANÁLISE DAS QUESTÕES OBJETIVAS Esta seção apresenta estatísticas para o conjunto de questões objetivas e uma análise gráfica de questões selecionadas dos componentes de Formação Geral ( 3.2.1 ), de Conhecimento Específico Bacharelado ( 3.2.2 ) e de Conhecimento Específico Licenciatura ( 3.2.3 ).
São também apresentadas e comparadas as médias das subpopulações caracterizadas por Grande Região. 3.2.1 Componente de Formação Geral A Tabela 3.7 apresenta as estatísticas básicas relativas às oito questões objetivas do componente da prova que abrange a Formação Geral dos estudantes.
A média do Brasil foi 53,1.
A menor média foi encontrada na região Norte ( 49,5 ), e a maior, na região Sudeste ( 56,8 ).
As demais médias foram : 52,4 na região Nordeste ; 52,4 na região Sul ; e 53,5 na região Centro-Oeste.
O desvio padrão do Brasil foi 21,0, sendo o maior desvio padrão encontrado na região Sudeste ( 21,8 ), e o menor, na região Nordeste ( 20,3 ).
Os demais desvios foram : 20,7 na região Norte ; 20,6 na região Sul ; e 21,2 na região Centro-Oeste.
As medianas ( 50,0 ), as notas máximas ( 100,0 ) e as notas mínimas ( 0,0 ) foram iguais para todas as regiões, exceto a mediana na região Sudeste ( 62,5 ).
A Tabela 3.8 apresenta o índice de facilidade e o índice de discriminação ( ponto bisserial ) para cada uma das questões objetivas do Componente de Formação Geral.
Quanto ao índice de facilidade, foram usadas as seguintes cores para diferenciar o nível de dificuldade da questão : Já quanto ao índice de discriminação, foram usadas as seguintes cores para qualificar a questão : As questões objetivas do Componente de Formação Geral, segundo o índice de facilidade, foram assim avaliadas : das oito questões, nenhuma teve o índice de facilidade classificado como muito fácil.
Duas questões foram tidas como fácil, por terem índice de acertos situado na faixa entre 0,61 e 0,85 ( de 61,0 \% a 85,0 \% de acertos ).
Cinco questões foram consideradas de dificuldade médio, situando -se no intervalo entre 0,41 e 0,60 do índice de facilidade, ou seja, houve entre 41,0 \% e 60,0 \% de acertos, enquanto uma questão foi classificada na categoria difícil, situando -se no intervalo entre 0,16 e 0,40.
Por fim, nenhuma das questões apresentou menos de 15 \% de acertos, razão pela qual seria classificada como muito difícil.
Como já comentado, para análise das questões objetivas, relativas à Formação Geral, segundo o poder de discriminação, utilizou -se o índice de discriminação ( ponto bisserial ).
Nesta análise, as questões foram assim avaliadas : seis das oito questões apresentaram índice acima ou igual a 0,40 e, assim, foram classificadas com índice muito bom para esse grupo de estudantes.
Duas questões tiveram índice de discriminação bom, com valor entre 0,30 e 0,39, para esse grupo de estudantes.
Nenhuma questão teve nível médio ou fraco de discriminação para esse grupo de estudantes.
O índice de facilidade variou de 0,25 a 0,81, e o de discriminação, de 0,36 a 0,53.
As seis questões com índice de discriminação muito bom figuraram entre os diversos níveis de dificuldade desse conjunto : duas classificadas na categoria fácil ( questões 4 e 8 ) do índice de facilidade, e as outras quatro, na categoria médio ( questões 1, 3, 6 e 7 ).
Em particular, a questão 7 foi a que apresentou o maior poder discriminatório, com índice 0,53, porém foi considerada média em termos de facilidade, com uma proporção de 0,43 acertos.
O máximo de acertos foi alcançado pela questão 8 com um índice de facilidade de 0,81.
A questão de número 1 apresentou um índice de facilidade de 0,50, ou seja, a metade dos estudantes conseguiu resolvê -la, dentro do universo de participantes.
Seu índice de discriminação foi muito bom ( 0,40 ).
Já a questão 5 obteve índice de discriminação bom, 0,36, porém seu índice de facilidade foi difícil ( 0,25 ).
O Gráfico 3.4, para exemplificar, analisa o comportamento da questão de número 7 de Formação Geral.
Trata -se de uma questão considerada mediana, com relação a facilidade e a que obteve o maior índice de discriminação dessa parte da prova.
Neste gráfico, cada uma das cinco curvas representa o percentual de respostas em determinada alternativa da questão, em função do número de acertos dos estudantes nessa parte da prova ( Formação Geral/Múltipla Escolha ), antes de possíveis eliminações pelo critério do ponto bisserial.
A curva em vermelho corresponde à alternativa E, a correta para esta questão.
Assim, observa -se que entre os estudantes com menor número de acertos, nessa parte do exame, a situação mais frequente foi a escolha de uma das alternativas incorretas : a alternativa C ( em preto ) ou D ( em roxo ).
À medida em que o número de acertos aumenta, indicando desempenho melhor nessa parte da prova, aumenta concomitantemente a proporção de estudantes que selecionaram a alternativa correta E, atingindo 100 \% para os estudantes com 8 acertos.
Essa análise permite verificar como a questão discriminou os grupos de desempenho, justificando o alto índice obtido na questão.
Os gráficos relativos às demais questões de Formação Geral constam do Anexo I. 3.2.2 Componente de Conhecimento Específico – Letras - Português ( Bacharelado ) A Tabela 3.9 apresenta as estatísticas básicas em relação às questões objetivas do Componente de Conhecimento Específico da prova de Letras - Português ( Bacharelado ), por Grande Região.
A média do Brasil deste componente foi de 48,4.
A menor média foi observada na região Norte ( 42,9 ), e a maior, na região Sudeste ( 49,7 ).
O desvio padrão de todo o Brasil foi 15,8, sendo o menor desvio padrão encontrado na região Nordeste ( 12,9 ), e o maior, na região Sul ( 19,1 ).
A mediana de todo o Brasil foi 48,0, a mesma encontrada na região Sul.
As regiões Norte ( 40,0 ), Nordeste ( 46,0 ) e Centro-Oeste ( 44,0 ) apresentaram valores menores para a mediana, enquanto a região Sudeste ( 52,0 ) apresentou um valor maior.
A nota máxima da prova foi 92,0, obtida nas questões objetivas do Componente de Conhecimento Específico, por, pelo menos, um estudante na região Centro-Oeste, enquanto nas demais regiões, a nota máxima da prova foi : 72,0 na região Norte ; 84,0 nas regiões Nordeste e Sudeste ; e 88,0 na região Sul.
A nota mínima foi zero nas regiões Sudeste e Sul, enquanto nas demais regiões foi : 28,0 na região Norte ; 24,0 na região Nordeste ; e 12,0 na região Centro-Oeste.
A Tabela 3.10 apresenta os índices de facilidade e discriminação ( ponto bisserial ) das questões objetivas do Componente de Conhecimento Específico para os estudantes de Letras - Português ( Bacharelado ).
Para facilitar a diferenciação das questões usaram -se as mesmas cores da Tabela 3.8 para as diferentes classificações dos índices de facilidade e de discriminação.
Dentre as questões objetivas da parte da prova relativa ao Componente de Conhecimento Específico, nenhuma das questões foi anulada pela Comissão.
Desse modo, a classificação, quanto ao índice de facilidade, foi estabelecida com base em todas as 27 questões.
A partir dos índices obtidos, pode -se concluir que um pouco mais de um terço das questões objetivas da prova foi considerado, pelo menos, difícil : das 27 questões, 11 foram classificadas como difícil ( a classificação modal ).
Seis questões foram classificadas como fácil, e outras dez consideradas como médio.
Já quanto aos índices de discriminação das questões objetivas do Componente de Conhecimento Específico da prova, tem -se como resultado a seguinte classificação : 13 das 27 questões foram consideradas como boas, enquanto quatro delas tiveram índice de discriminação muito bom.
Assim, para a maioria das questões – 17 em 27 – os índices de discriminação foram bom ou muito bom.
Dentre as demais, oito delas foram classificadas como médio e outras duas como fraco, sendo dez, por conseguinte, a quantidade de questões nos dois patamares mais baixos de discriminação.
Constata -se, assim, que a prova – no que se refere ao Componente de Conhecimento Específico – possuía boa capacidade de discriminar entre aqueles que dominam ou não o conteúdo.
Dentre as questões que alcançaram os maiores índices de discriminação, quatro delas, as de números 17, 27, 30 e 33, foram classificadas com índice de discriminação muito bom, situando -se no intervalo de 0,40 a 0,47 do índice, três delas ( questões 17, 27 e 30 ) foram classificadas na categoria fácil, quanto ao índice de facilidade, e a outra ( questão 33 ), na categoria médio.
A questão de número 21 foi a A Tabela 3.10 apresenta os índices de facilidade e discriminação ( ponto bisserial ) das questões objetivas do Componente de Conhecimento Específico para os estudantes de Letras - Português ( Bacharelado ).
Para facilitar a diferenciação das questões usaram -se as mesmas cores da Tabela 3.8 para as diferentes classificações dos índices de facilidade e de discriminação.
Dentre as questões objetivas da parte da prova relativa ao Componente de Conhecimento Específico, nenhuma das questões foi anulada pela Comissão.
Desse modo, a classificação, quanto ao índice de facilidade, foi estabelecida com base em todas as 27 questões.
A partir dos índices obtidos, pode -se concluir que um pouco mais de um terço das questões objetivas da prova foi considerado, pelo menos, difícil : das 27 questões, 11 foram classificadas como difícil ( a classificação modal ).
Seis questões foram classificadas como fácil, e outras dez consideradas como médio.
Já quanto aos índices de discriminação das questões objetivas do Componente de Conhecimento Específico da prova, tem -se como resultado a seguinte classificação : 13 das 27 questões foram consideradas como boas, enquanto quatro delas tiveram índice de discriminação muito bom.
Assim, para a maioria das questões – 17 em 27 – os índices de discriminação foram bom ou muito bom.
Dentre as demais, oito delas foram classificadas como médio e outras duas como fraco, sendo dez, por conseguinte, a quantidade de questões nos dois patamares mais baixos de discriminação.
Constata -se, assim, que a prova – no que se refere ao Componente de Conhecimento Específico – possuía boa capacidade de discriminar entre aqueles que dominam ou não o conteúdo.
Dentre as questões que alcançaram os maiores índices de discriminação, quatro delas, as de números 17, 27, 30 e 33, foram classificadas com índice de discriminação muito bom, situando -se no intervalo de 0,40 a 0,47 do índice, três delas ( questões 17, 27 e 30 ) foram classificadas na categoria fácil, quanto ao índice de facilidade, e a outra ( questão 33 ), na categoria médio.
A questão de número 21 foi a A título de exemplo das análises do comportamento das questões objetivas, o Gráfico 3.5 analisa a questão 30 do Componente de Conhecimento Específico.
Esta mais difícil dentre as 27 questões específicas, com baixo índice de facilidade, apenas 19,0 \% de acertos.
Essa questão apresentou poder discriminatório igualmente baixo, 0,09, o que comprova ter sido esta questão a mais difícil para os estudantes.
Destaca -se, também, a questão 23, com índice de facilidade 0,28, o que, em termos percentuais, corresponde a 28,0 \% de estudantes que responderam acertadamente.
Já, 0,19 foi o seu índice de discriminação.
Tais questões foram, portanto, pelo critério ponto bisserial, consideradas inadequadas.
Por isso, as questões 21 e 23 foram as únicas eliminadas do cômputo da nota final.
A título de exemplo das análises do comportamento das questões objetivas, o Gráfico 3.5 analisa a questão 30 do Componente de Conhecimento Específico.
Esta questão foi classificada como fácil.
Apresentou índice de facilidade 0,62, ou seja, 62,0 \% dos estudantes assinalaram acertadamente a opção E, correspondente ao gabarito.
Seu índice de discriminação foi igual a 0,47, classificado como muito bom, também sendo esta questão a que apresentou o maior índice discriminatório.
Neste gráfico, cada uma das cinco curvas representa o percentual de respostas em determinada alternativa da questão 30, em função do número de acertos dos estudantes nessa parte da prova, antes de possíveis eliminações de questões pelo critério do ponto bisserial.
A alternativa correta C, representada no gráfico pela curva em preto, foi escolhida em maiores proporções pelos estudantes com desempenho melhor nessa parte da prova.
Já as alternativas incorretas, também denominadas distratores, foram selecionadas, principalmente, por aqueles com notas mais baixas.
Neste caso também a soma não é sempre 100 \% por causa das questões não respondidas ou com mais de uma opção marcada.
Aqueles com nota zero, na sua quase totalidade deixaram esta questão em branco ou marcaram mais de uma alternativa, comportamento considerado inválido.
A proporção de estudantes que selecionou a resposta correta C aumenta gradativamente, chegando a atingir 100 \% para 21 acertos ou mais, enquanto a proporção dos que escolheram alternativas incorretas decai, a partir de três acertos, como função do número de acertos nessa parte da prova.
Os gráficos relativos às demais questões do Componente de Conhecimento Específico constam do Anexo I. 3.2.3 Componente de Conhecimento Específico – Letras - Português ( Licenciatura ) A Tabela 3.11 apresenta as estatísticas básicas em relação às questões objetivas do Componente de Conhecimento Específico da prova de Letras - Português ( Licenciatura ), por Grande Região.
A média do Brasil deste componente foi de 40,5.
A menor média foi observada na região Norte ( 37,7 ), e a maior, na região Sudeste ( 43,8 ).
O desvio padrão de todo o Brasil foi 15,6, sendo o menor desvio padrão encontrado na região Norte ( 14,5 ), e o maior, na região Sudeste ( 16,9 ).
A mediana de todo o Brasil foi 40,9, a mesma encontrada na região Nordeste.
As regiões Norte, Sul e Centro-Oeste ( 36,4 em cada ) apresentaram valores menores para a mediana, enquanto a região Sudeste ( 45,5 ) apresentou valor.
A nota máxima da prova foi 95,5, obtida nas questões objetivas do Componente de Conhecimento Específico, por, pelo menos, um estudante nas regiões Nordeste e Sudeste, enquanto nas demais regiões, a nota máxima da prova foi : 90,9 na região Norte ; 86,4 na região Sul ; e 81,8 na região Centro-Oeste.
Em todas as regiões a nota mínima foi zero, sem exceção.
A Tabela 3.12 apresenta os índices de facilidade e discriminação ( ponto bisserial ) das questões objetivas do Componente de Conhecimento Específico para os estudantes de Letras - Português ( Licenciatura ).
Para facilitar a diferenciação das questões usaram -se as mesmas cores das Tabelas 3.8 e 3.10 para as diferentes classificações dos índices de facilidade e de discriminação.
Dentre as questões objetivas da parte da prova relativa ao Componente de Conhecimento Específico, nenhuma das questões foi anulada pela Comissão.
Desse modo, a classificação, quanto ao índice de facilidade, foi estabelecida com base em todas 27 questões.
A partir dos índices obtidos, pode -se concluir que mais de metade das questões objetivas da prova foi considerado, pelo menos, difícil : das 27 questões, 16 foram classificadas como difícil ( a classificação modal ) ou como muito difícil.
Uma questão foi classificada como fácil ou como muito fácil, e outras dez consideradas como médio.
Já quanto aos índices de discriminação das questões objetivas do Componente de Conhecimento Específico da prova, tem -se como resultado a seguinte classificação : seis das 27 questões foram consideradas como boas, enquanto outras seis delas tiveram índice de discriminação muito bom.
Assim, para um pouco menos de metade das questões – 12 em 27 – os índices de discriminação foram bom ou muito bom.
Dentre as demais, dez delas foram classificadas como médio e outras cinco como fraco, sendo 15, por conseguinte, a quantidade de questões nos dois patamares mais baixos de discriminação.
Constata -se, assim, que a prova – no que se refere ao Componente de Conhecimento Específico – possuía capacidade mediana de discriminar entre aqueles que dominam ou não o conteúdo.
Dentre as questões que alcançaram os maiores índices de discriminação, seis delas, as de números 13, 14, 17, 28, 30 e 31, foram classificadas com índice de discriminação muito bom, situando -se no intervalo de 0,40 a 0,45 do índice, uma delas ( questão 13 ) foram classificadas na categoria fácil, quanto ao índice de facilidade, e as outras cinco ( questões 14, 17, 28, 30 e 31 ), na categoria médio.
A questão de número 27 foi a mais difícil dentre as 27 questões específicas, com baixo índice de facilidade, apenas 14,0 \% de acertos.
Essa questão apresentou poder discriminatório igualmente baixo, 0,12, o que comprova ter sido esta questão a mais difícil para os estudantes.
Destaca -se, também, a questão 33, com índice de facilidade 0,16, o que, em termos percentuais, corresponde a 16,0 \% de estudantes que responderam acertadamente.
Já, 0,11 foi o seu índice de discriminação.
Tais questões foram, portanto, pelo critério ponto bisserial, consideradas inadequadas.
Por isso, as questões 27 e 33 foram eliminadas do cômputo da nota final.
Além destas duas, as demais questões com índice de discriminação fraco, questões 21, 22 e 23 também não foram utilizadas no cômputo final das notas, num total de cinco questões eliminadas.
A título de exemplo das análises do comportamento das questões objetivas, o Gráfico 3.6 analisa a questão 17 do Componente de Conhecimento Específico.
Esta questão foi classificada como médio.
Apresentou índice de facilidade 0,55, ou seja, 55,0 \% dos estudantes assinalaram acertadamente a opção B, correspondente ao gabarito.
Seu índice de discriminação foi igual a 0,45, classificado como muito bom, também sendo esta questão a que apresentou o maior índice discriminatório.
Neste gráfico, cada uma das cinco curvas representa o percentual de respostas em determinada alternativa da questão 17, em função do número de acertos dos estudantes nessa parte da prova, antes de possíveis eliminações de questões pelo critério do ponto bisserial.
A alternativa correta B, representada no gráfico pela curva em verde, foi escolhida em maiores proporções pelos estudantes com desempenho melhor nessa parte da prova.
Já as alternativas incorretas, também denominadas distratores, foram selecionadas, principalmente, por aqueles com notas mais baixas.
Neste caso também a soma não é sempre 100 \% por causa das questões não respondidas ou com mais de uma opção marcada.
Aqueles com nota zero, na sua quase totalidade deixaram esta questão em branco ou marcaram mais de uma alternativa, comportamento considerado inválido.
A proporção de estudantes que selecionou a resposta correta B aumenta gradativamente, chegando a atingir 100 \% para 21 acertos ou mais, enquanto a proporção dos que escolheram alternativas incorretas decai, a partir do primeiro acerto, como função do número de acertos nessa parte da prova.
Os gráficos relativos às demais questões do Componente de Conhecimento Específico constam do Anexo I. 3.3 ANÁLISE DAS QUESTÕES DISCURSIVAS Esta seção apresenta estatísticas selecionadas e histogramas das Questões Discursivas de Formação Geral ( 3.3.1 ) e Conhecimento Específico ( 3.3.2 ).
São também apresentadas e comparadas as médias de subpopulações caracterizadas por Grande Região. 3.3.1 Componente de Formação Geral As análises dos resultados de desempenho dos estudantes de Letras - Português ( Bacharelado e Licenciatura ) nas duas questões discursivas relativas à Formação Geral, considerando -se as notas de conteúdo e de desempenho linguístico, encontram -se na Tabela 3.13 e no Gráfico 3.7.
Na Tabela 3.13, observa -se que a nota média nesse conjunto de questões foi próxima da obtida nas objetivas.
Os estudantes, de todo o Brasil, obtiveram, em Formação Geral, média 53,1 nas questões objetivas e 52,6 nas questões discursivas.
No entanto, pode -se notar um aumento do desvio padrão de 21,0 nas questões objetivas do Componente de Formação Geral dos estudantes de todo o Brasil, para 25,0 nas questões discursivas do mesmo componente.
A maior média foi obtida na região Sudeste ( 54,7 ), e a menor, na região Norte ( 48,3 ).
A mediana de todo o Brasil, neste componente, foi 57,0, a mesma obtida na região Nordeste.
Na região Sudeste a mediana foi maior, 61,0.
Nas demais regiões foi menor : Norte ( 52,5 ), Sul ( 56,5 ) e Centro-Oeste ( 56,0 ).
A nota máxima ( 98,5 ) foram as mesmas nas regiões Sudeste e Centro-Oeste ; sendo 97,5 na região Norte e 98,0 nas regiões Nordeste e Sul.
A nota mínima ( 0,0 ) foram as mesmas em todas as regiões do Brasil, sem exceção.
O Gráfico 3.7 representa a distribuição das notas nas questões discursivas do Componente de Formação Geral, segundo a opção de Bacharelado/Licenciatura.
As modas destas distribuições ocorrem no intervalo ( 70 ; 80 ] para o Bacharelado e no intervalo ( 60:70 ] para Licenciatura.
Para as duas habilitações observa -se outro máximo local no intervalo [ 0 ; 10 ], sendo que nesse intervalo incluem -se além das notas zero, a frequência de estudantes que deixaram este tipo de questão em branco.
Na sequência, os resultados verificados para cada uma das questões discursivas de Formação Geral serão apresentados, estabelecendo -se relações com os conteúdos abordados em cada uma delas.
Os comentários da Banca de docentes corretores, a respeito do observado na correção das respostas dos estudantes, suas impressões e conclusões serão apresentados junto à análise de cada questão.
Cumpre esclarecer que, tendo em vista que as questões discursivas de Formação Geral são padronizadas, ou seja, constam de todas as provas, os comentários da Banca são os mesmos para todas as carreiras acadêmicas, sendo direcionados a todos os estudantes que participaram do ENADE/2014.
A seguir, serão analisados os desempenhos da Área de Letras - Português ( Bacharelado e Licenciatura ) nas duas questões discursivas de Formação Geral do ENADE/2014, comparando -se os resultados obtidos com comentários para cada questão. 3.3.1.1 Análise de conteúdo da Questão Discursiva 1 do Componente de Formação Geral Os dados de Letras - Português ( Bacharelado e Licenciatura ), obtidos a partir das respostas à questão 1, encontram -se na Tabela 3.14 e no Gráfico 3.8.
Nessa questão – de pior desempenho dentre as duas de Formação Geral – os estudantes, de todo o Brasil, tiveram média 49,3.
A maior média para a questão 1 foi obtida na região Sudeste ( 53,5 ), e a menor, na região Norte ( 43,2 ).
Quanto à variabilidade das notas, o desvio padrão de todo o Brasil foi 31,9.
O menor desvio padrão foi obtido na região Nordeste ( 31,0 ), e o maior desvio padrão foi obtido na região Sudeste ( 33,0 ).
As medianas Brasil, como um todo, e de todas as regiões foram iguais ( 50,0 ), exceto na região Sudeste ( 60,0 ).
As notas máximas e mínimas da questão discursiva 1 foram as mesmas para todas as regiões do Brasil, respectivamente, 100,0 e 0,0.
O Gráfico 3.8 mostra a distribuição das notas na questão discursiva 1 do Componente de Formação Geral para estudantes do Bacharelado e da Licenciatura.
Observa -se que a maior frequência corresponde aos estudantes que obtiveram notas no intervalo ( 70 ; 80 ], tanto para os de Bacharelado quanto para os de Licenciatura. 3.3.1.2 Comentários sobre a correção de Conteúdo das respostas à Questão Discursiva 1 O enunciado apresentado da Questão 1 era claro e foi redigido em linguagem acessível.
A questão abordava o tema da mobilidade urbana, relacionando -a aos ideais de desenvolvimento urbano sustentável estabelecidos pela ONU, do ponto de vista ecológico, cultural, político, institucional, social e econômico.
O enunciado apresentava uma tabela, sem referência à fonte de publicação, que detalhava percentualmente as modalidades de deslocamento da população urbana brasileira, com base em duas distinções : motorizado versus não motorizado ; coletivo versus individual.
O comando da questão solicitava, então, que o estudante desenvolvesse um texto dissertativo composto de duas partes : análise das consequências do transporte motorizado para o desenvolvimento sustentável e elaboração de duas propostas de intervenção pautadas pelo incentivo ao uso da bicicleta, meio de transporte não motorizado e, portanto, não poluente.
A temática dos efeitos da industrialização, com foco na produção de automóveis, e consequências para o meio ambiente é atual.
Além disso, a questão provoca a reflexão sobre a necessidade de adequação das cidades ao número de habitantes, no contexto da utilização dos meios de transportes de massa : metroviário, ferroviário, por ônibus.
Esperava -se, assim, que o concluinte recorresse a diferentes argumentos para analisar as consequências do transporte motorizado, explicitados pelo texto motivador, e propusesse duas ações de intervenção, as quais considerasse relevantes para o incremento ao uso de bicicleta, que, preferencialmente, deveriam manter coerência com a argumentação desenvolvida na primeira parte da resposta.
O conflito entre a busca pelo transporte individual e o coletivo é inerente a todas as grandes cidades do mundo.
No Brasil, esse conflito emerge com maior gravidade, como efeito do aumento da quantidade de veículos automotores ( automóveis e motocicletas ) e da baixa qualidade dos transportes coletivos.
Tal associação incentiva o uso de meios de transporte individual em detrimento da utilização dos coletivos.
O padrão de respostas era tecnicamente adequado ao enunciado e amplo o suficiente para abranger as variações de respostas apresentadas.
Deve ser ressaltado que a proteção ao ambiente ecologicamente equilibrado é tema conhecido dos concluintes dos cursos universitários, que deveriam possuir informações suficientes para desenvolver uma resposta abordando pelo menos dois dos aspectos listados no padrão de respostas como possibilidades de consequências ( item ‘ a ’ ) e duas propostas de ação de intervenção aceitáveis ( item ‘ b ’ ).
Segundo o padrão, concluintes poderiam desenvolver, quanto ao item ‘ a ’, os seguintes pontos : aumento da emissão de poluentes atmosféricos ; aumento da emissão de gases de efeito estufa ( CO2- dióxido de carbono, CO - monóxido de carbono, O3 - Ozônio ) ; aumento da poluição visual e sonora ; aumento da temperatura local e global ; aumento do consumo de combustíveis ; aumento de problemas de saúde ( cardíacos, respiratórios, dermatológicos ) ; aumento da frota de veículos, promovendo congestionamentos urbanos ; diminuição de áreas verdes ; desmatamento ; aumento das áreas de impermeabilização, resultando em enchentes, diminuição da infiltração da água e recarga de lençóis freáticos ; elevação dos custos de manutenção das cidades ( metroferrovias, rodovias, tratamento de água, limpeza da cidade, etc. ) ; necessidade de ampliação de vias trafegáveis ; necessidade de ampliação de áreas de estacionamento.
Em relação ao item ‘ b ’, os estudantes poderiam considerar os seguintes aspectos : construção de vias exclusivas para bicicletas ( ciclovias e ciclofaixas ) ; proposição de formas de integração entre o transporte por bicicletas, o metroviário e os ônibus coletivos, a fim de garantir segurança e conforto em momentos de adversidades climáticas e relevo acidentado ; pontos de aluguel e/ou empréstimos de bicicletas ; construção de bicicletários ; investimentos na segurança pública ; políticas de incentivo ao uso de bicicletas ( educação ambiental, qualidade de vida, saúde, propaganda ) ; implementação de políticas de crédito e de redução de custos das bicicletas.
Para o item ‘ a ’, foi encontrado um número expressivo de respostas adequadas.
As consequências mais citadas estavam relacionadas com a emissão de poluentes e o aumento dos engarrafamentos.
Houve críticas frequentes à qualidade do transporte público e à necessidade de locomoção rápida diante da distância das residências aos locais de trabalho.
Em tom de reivindicação, muitos estudantes defenderam solução para os problemas do transporte coletivo, quer metroviário, ferroviário ou por ônibus.
O uso da bicicleta também foi considerado por muitos como um meio de transporte benéfico à saúde, mas alguns apontaram a impossibilidade de sua utilização pelos idosos e deficientes.
No concernente ao item ‘ b ’, de modo quase uniforme citou -se a necessidade de ciclovias e ciclofaixas, bem como de se garantir segurança para os ciclistas.
A necessidade de manutenção das ciclovias e ciclofaixas existentes em algumas cidades foram também lembradas.
Com menos frequência também houve a indicação de maior conscientização da população, inclusive pela mídia, e pelo incentivo à aquisição desses veículos não motorizados, por meio da diminuição de impostos para redução do custo de aquisição.
Houve quem sugerisse, inclusive, uma política de doação de bicicletas.
Alguns mencionaram a necessidade de se assegurar integração das ciclovias com os demais modais, com a construção de bicicletários seguros diante das grandes distâncias a percorrer, notadamente nas metrópoles.
Também com base na dificuldade de o veículo não motorizado vencer as longas distâncias que o trabalhador deve percorrer nos centros urbanos para chegar ao local onde exerce a sua atividade produtiva, alguns estudantes argumentaram contra a utilização da bicicleta como forma de resolução da poluição ambiental.
O principal equívoco cometido na sugestão de ações ( item ‘ b ’ ) foi o de mencionar a necessidade de melhoria do meio ambiente dissociado do uso de bicicletas, como se requeria no comando da questão. bicicletas, como se requeria no comando da questão.
Nesse sentido, foram encontradas respostas que indicavam a caminhada como opção.
Foram consideradas respostas fracas ( notas de zero a 30 ) aquelas meramente opinativas, ou que apenas transcreveram elementos do enunciado, sem qualquer acréscimo ou apreciação crítica.
As respostas medianas ( notas entre 35 e 70 ) foram as que não apresentavam algum desenvolvimento para os tópicos citados ( duas consequências e duas ações de intervenção ) ou por falta de alguns deles.
Já as respostas boas ( notas de 75 a 100 ) continham os tópicos considerados pelo padrão de respostas, com algum desenvolvimento, pelo menos, regular de argumentação.
As respostas dos estudantes revelaram um adequado conhecimento quanto aos efeitos do transporte motorizado para o ambiente e que a temática da ecologia parece ser tema bem compreendido, aparecendo nas respostas como uma preocupação para as futuras gerações.
Nesse sentido, algumas políticas públicas foram mencionadas apropriadamente para assegurar o desenvolvimento sustentável pelos respondentes, tais como : educação, habitação, saneamento, saúde, transporte e segurança.
No entanto, de modo geral, as respostas demostraram que os estudantes têm muita dificuldade na expressão escrita do pensamento, como se pode constatar pela avaliação do desempenho linguístico que ficou a cargo de uma banca específica, formada por profissionais da área de Língua Portuguesa. 3.3.1.3 Análise de Conteúdo da Questão Discursiva 2 do Componente de Formação Geral A Tabela 3.15 mostra que o desempenho médio dos estudantes na questão discursiva 2 ( média 53,3 ) foi superior ao obtido na questão discursiva 1 ( média 49,3 ).
A região Nordeste foi aquela cuja média, nessa questão, foi maior ( 54,8 ), e a de menor média foi a região Centro-Oeste ( 49,6 ).
Quanto à variabilidade das notas, o desvio padrão de todo o Brasil foi 32,0, ligeiramente superior ao obtido na questão discursiva 1 ( 31,9 ).
O maior desvio nessa questão foi obtido na região Centro-Oeste ( 33,5 ), enquanto o menor foi obtido na região Sul ( 30,2 ).
As medianas ( 60,0 ), as notas máximas ( 100,0 ) e as notas mínimas ( 0,0 ) foram iguais para todas as regiões, sem exceção.
O Gráfico 3.9 mostra a distribuição das notas na questão discursiva 2 do Componente de Formação Geral para estudantes do Bacharelado e da Licenciatura.
Observa -se que a maior frequência corresponde aos estudantes que obtiveram nota no intervalo ( 70 ; 80 ] para o Bacharelado, e no intervalo ( 60 ; 70 ] para a Licenciatura. 3.3.1.4 Comentários sobre a correção de Conteúdo das respostas à Questão Discursiva 2 A questão 2 abordava o tema da violência urbana a partir de um texto motivador recolhido em um site de um jornal brasileiro.
O fragmento relatava a prisão de três jovens por agressão a um outro jovem, menor de idade, guardador de carros.
A notícia trazia informações sobre o estado do jovem agredido e sobre as possíveis penas a serem aplicadas aos agressores.
O enunciado solicitava, então, que o estudante desenvolvesse um texto dissertativo composto de duas partes : análise de duas causas do tipo de violência descrito no texto e apresentação de dois fatores que contribuiriam para evitar a agressão relatada.
Como o texto motivador reduzia -se ao relato de um episódio de violência urbana, o enunciado não oferecia argumentos que ajudassem o estudante a elaborar uma análise das possíveis causas do fato.
As propostas a serem apresentadas, no sentido de se evitar em situações como a relatada, deveriam manter coerência com a análise das causas, desenvolvida na primeira parte da resposta.
O enunciado era claro e a matéria permitia identificar tratar -se de um conflito tipicamente urbano no qual se envolveram jovens de baixa renda em disputa por um território onde exerciam trabalho informal.
Sendo essa uma das hipóteses albergada sobre o amplo tema da violência urbana no qual se desenvolve o relato.
A temática da questão é bastante conhecida, não somente por debates no âmbito universitário, especialmente na área de Ciências Sociais, bem como pela ênfase que é dada pela mídia onde assume contornos, muitas vezes, maiores do que o problema em si.
Assim, por tratar -se de tema bastante presente e debatido no cotidiano das pessoas, a questão pode ser considerada de baixa complexidade.
O padrão de resposta apresentava uma gama de possibilidades de causas da violência aceitas como corretas para o item ‘ a ’, tais como : problemas relacionados à educação ( baixa escolaridade, evasão escolar, qualidade da educação, distanciamento entre escola e a realidade social, tempo de permanência na escola ) ; desigualdades socioculturais ( gênero, etnia, economia, etc. ) ; desemprego e falta de qualificação profissional ; precariedade da segurança pública ; uso de drogas ; desvalorização da vida humana ; banalização da violência ; sensação de impunidade ; ausência de políticas sociais ; degradação da vida urbana ; desconhecimento ou desrespeito aos direitos humanos e constitucionais ; desestruturação familiar ; desvalorização de princípios éticos e morais.
No concernente ao item ‘ b ’, o padrão também previa a possibilidade de apresentação de diversos fatores que podem contribuir para evitar a violência : políticas de segurança mais efetivas ; políticas públicas de melhoria das condições socioeconômicas ; maior consciência cidadã e respeito à vida ; melhor distribuição de renda ; melhoria da educação ( aumento da escolaridade, redução da evasão escolar, qualidade da educação, aproximação entre a escola e a realidade social, aumento do tempo de permanência na escola ) ; aumento da oferta de emprego e melhoria da qualificação profissional ; medidas preventivas ao uso de drogas ; maior eficácia do sistema judiciário ; revisão da legislação penal ; valorização de princípios éticos, morais e familiares.
O enunciado da questão permitia que os estudantes, em suas respostas, pautassem o problema por diversos ângulos.
Dessa forma, o padrão buscou indicar diversas possibilidades a serem identificadas pelo corretor na leitura da resposta do estudante.
As notas foram atribuídas considerando o desenvolvimento da resposta e a indicação de ao menos dois dos itens exigidos pelo comando da questão nos itens ‘ a ’ ( duas causas ) e ‘ b ’ ( dois fatores ), dentre os considerados pelo padrão de respostas.
No entanto, o padrão exigia que as respostas estivessem afinadas com o perfil profissional previsto no Art. 3º da Portaria nº 255, de 02 de junho de 2014, onde se lê : Art. 3º No componente de Formação Geral serão considerados os seguintes elementos integrantes do perfil profissional : atitude ética ; comprometimento social ; compreensão de temas que transcendam ao ambiente próprio de sua formação, relevantes para a realidade social ; espírito científico, humanístico e reflexivo ; capacidade de análise crítica e integradora da realidade ; e aptidão para socializar conhecimentos em vários contextos e públicos diferenciados.
Nesse sentido, teses repressivas não conformadas com o texto constitucional e amparados pelos direitos humanos, não foram consideradas.
Questões como a majoração das penas, instituição de penas mais duras e criminalização da conduta dos menores desviados, só foram consideradas no caso de serem tratadas como menção a um debate que está presente, nunca como medida a ser aplicada de forma arbitrária e inconstitucional.
Também não foram previstas visões mais autoritárias e repressoras como a defesa de proibição do trabalho informal, como é o caso dos “ flanelinhas ”, e a defesa de não fornecer dinheiro aos “ guardadores ” de carro, identificado como esmolas por alguns.
Não foram consideradas respostas que meramente reproduziam o enunciado sem qualquer reflexão ou análise.
Aquelas que utilizaram os elementos do enunciado, mesmo que apenas transcritos, mas desenvolveram algum raciocínio pessoal ou indicaram causas e/ou fatores foram valoradas.
As causas apontadas com maior frequência para a violência foram : a ausência de educação e o desemprego, com soluções correlatas de educação profissional e de qualidade e abertura de postos de trabalho.
A questão das drogas também foi bastante mencionada, e a solução correlata mais indicada foi o acolhimento de usuários de drogas.
Outro tema bastante citado foi a dissolução das famílias e dos valores familiares.
Tendo em vista que o texto falava de moradores de rua, a ausência de moradia apareceu em muitas respostas, bem como a disputa territorial e a necessidade de sobrevivência.
Nesse sentido, as propostas apontavam para a necessidade de ajuda da área de Serviço Social das prefeituras e atendimento psicológico gratuito.
Numa ótica mais repressora, surgiram críticas quanto à leniência das leis, acompanhadas da indicação da necessidade de atuação mais rigorosa do Poder Judiciário, como também da aplicação de penas maiores, negativa de esmolas, proibição do trabalho informal ( guardador de carros ), dentre outras.
Algumas respostas indicaram como causa o modelo capitalista, o consumismo e a ganância.
Outras indicaram o trabalho infantil e a correlata necessidade de aplicação do ECA ( Estatuto da Criança e do Adolescente ).
Em relação ao tema violência, inúmeras respostas apresentaram visões de intolerância quanto aos menores infratores, postulando a diminuição da idade para efeitos de criminalização, bem como defendendo a majoração de penas.
Essa perspectiva revela uma cisão social, de certa forma estimulada pelo tipo de divulgação da mídia quanto aos temas em foco, o que, sem dúvida, ficou evidenciado em algumas respostas.
A existência de inúmeras respostas pugnando pela maior repressão estatal e pelo aumento das penas, acrescido da busca pela diminuição da idade penal, revela que muitos estudantes não foram atingidos por uma formação superior que se exige humanista, ética e comprometida socialmente, que busque os fundamentos para resolver o problema da violência por meio de práticas democráticas e de inclusão social.
As respostas desse tipo também revelaram um alto grau de intolerância para com os moradores de rua, associando -os muitas vezes ao tráfico de drogas, à exploração indevida do espaço público, bem como a crimes como extorsão dos motoristas quando pedem dinheiro para guardar os veículos.
Nessa trilha, verifica -se lacuna na discussão de temas importantes para a formação de profissionais de nível superior tais como : sociodiversidade, multiculturalismo e violência ; tolerância/intolerância ; inclusão/exclusão.
Para aqueles que responderam em conformidade com o padrão de resposta verificou -se uma plena formação social, conclamando pelo auxílio não só do Estado, mas também da Sociedade no amparo daqueles que não têm as mínimas condições de sobrevivência e que estão nas ruas. 3.3.1.5 Análise de Língua Portuguesa das Questões Discursivas do Componente de Formação Geral Os dados de Letras - Português, obtidos a partir das respostas às questões discursivas do Componente de Formação Geral no que tange à Língua Portuguesa, encontram -se na Tabela 3.16 e no Gráfico 3.10.
Nesse aspecto, os estudantes, de todo o Brasil, tiveram média 58,0.
A maior média com respeito à Língua Portuguesa foi obtida na região Sul ( 59,6 ), e a menor, na região Norte ( 54,7 ).
Quanto à variabilidade das notas, o desvio padrão de todo o Brasil foi 23,2.
O menor desvio padrão foi obtido na região Sul ( 20,7 ) e o maior desvio padrão foi obtido na região Centro-Oeste ( 25,5 ).
A mediana das notas de Língua Portuguesa foi 65,0 para quase todas as regiões do Brasil, menos para as regiões Norte ( 62,5 ) e Sudeste ( 67,5 ).
A nota máxima para todo o Brasil foi de 95,0, com, pelo menos, um estudante tirando essa nota na região Norte, enquanto nas demais regiões a nota máxima foi 92,5 nas regiões Nordeste, Sudeste e Centro-Oeste ; e 90,0 na região Sul.
Além disso, a nota mínima foi zero em todas as regiões do país, sem exceção.
O Gráfico 3.10 mostra a distribuição das notas de Língua Portuguesa do Componente de Formação Geral.
Observa -se que a maior frequência corresponde aos estudantes que obtiveram nota no intervalo ( 70 ; 80 ] para o Bacharelado, e no intervalo ( 60 ; 70 ] para a Licenciatura. 3.3.1.6 Comentários sobre a correção das respostas de Formação Geral com respeito à Língua Portuguesa Ao encaminhar as questões 1 e 2 na direção da produção de um texto dissertativo, esperava -se que o estudante utilizasse seus conhecimentos sobre o assunto e estruturasse seus textos de acordo com as exigências do registro formal próprio dessa situação comunicativa.
Essa configuração determina exigências quanto : à adequação da seleção vocabular, ao desenvolvimento do conteúdo, à estruturação sintática dos períodos, à organização lógica das ideias, à utilização de procedimentos de encadeamento textual e referenciação, à obediência às exigências morfossintáticas próprias da modalidade escrita da norma-padrão, ao respeito às regras ortográficas e às regras de acentuação gráfica.
O padrão de respostas utilizado na avaliação considerou os aspectos relevantes ao bom desempenho linguístico como competências distintas, de modo a permitir um mapeamento detalhado do domínio dos recursos disponíveis na Língua Portuguesa para a comunicação escrita formal.
Com base nesse objetivo, foram avaliados os seguintes aspectos : a ) Estruturação textual condizente com o gênero solicitado e o modo de organização textual expositivo adequado ao gênero – essa competência envolve : a estruturação sintática condizente com o padrão da modalidade escrita formal da língua portuguesa, de modo a garantir a clareza necessária ; a distribuição do conteúdo do texto em parágrafos, de modo a garantir a sua organização temática ; a utilização de operadores discursivos que contribuam para a progressão temática do texto, estabelecendo relações lógicas entre as ideias apresentadas, tanto do ponto de vista intrafrasal, como do interfrasal ; a utilização de procedimentos de referenciação lexical e pronominal que permitam a retomada de referentes textuais ; o respeito às regras de pontuação como fator de estruturação do período.
Espera -se, portanto, que o estudante recorra a procedimentos linguístico-discursivos para organizar seu texto, permitindo o encadeamento lógico entre suas partes de forma a garantir a progressão e a coerência textuais.
Isso significa que os seguintes procedimentos foram penalizados, de acordo com o padrão de respostas proposto :  a estrutura lógico-gramatical do texto fica comprometida pela elaboração de frases fragmentadas ;  sequência justaposta de ideias sem encaixamentos sintáticos, reproduzindo hábitos da oralidade ;  elaboração de frase com apenas oração subordinada, sem oração principal ;  emprego equivocado do conector ( preposição, conjunção, pronome relativo, alguns advérbios e locuções adverbiais ) que não estabeleça relação lógica entre dois trechos do texto e prejudique a compreensão da mensagem ;  emprego do pronome relativo sem a preposição, quando obrigatória ;  repetição ou substituição inadequada de palavras sem utilização dos recursos oferecidos pela língua ( pronome, advérbio, artigo, sinônimo ) ;  emprego inadequado dos pronomes relativos “ cujo(a) ” e “ onde ” ;  utilização inadequada dos sinais de pontuação que comprometa a clareza textual. b ) Respeito às convenções ortográficas da norma-padrão da Língua Portuguesa – essa competência envolve o domínio das regras de acentuação gráfica e da grafia padrão das palavras ( com ausência de abreviaturas próprias da linguagem da internet ), de acordo com as convenções estabelecidas pela legislação em vigor e consubstanciadas no Vocabulário Ortográfico da Língua Portuguesa, editado pela Academia Brasileira de Letras ( com aceitação da legislação anterior, no caso das regras relativas ao uso do hífen e da acentuação gráfica ).
Espera -se que o participante :  grafe corretamente as palavras ;  respeite as regras de acentuação gráfica ;  empregue maiúsculas em início de frase, em nomes próprios de pessoas, lugares ou instituições ;  evite abreviações como p/, vc, tb, pq, tá, né, usadas muitas vezes em escrita informal e na internet ;  obedeça às regras de separação de sílabas no final da linha. c ) Domínio dos diferentes aspectos morfossintáticos próprios da modalidade escrita formal da norma-padrão da Língua Portuguesa – essa competência envolve : a concordância nominal, a concordância verbal, a regência nominal, a regência verbal, a flexão nominal, a flexão verbal, a correlação entre os tempos verbais, a colocação pronominal e a utilização de sinais de pontuação que contribuam para a organização lógica da frase e do texto.
Espera -se que o participante :  flexione o verbo para estabelecer concordância de número com o sujeito da frase ;  flexione o artigo, o adjetivo e o pronome para concordar em número e em gênero com o substantivo a que se referem ;  observe a regência nominal e a verbal, utilizando a preposição adequada depois de um substantivo, um verbo ou um adjetivo ;  empregue adequadamente o acento grave indicador de crase nos casos em que se fizer necessário ;  obedeça às regras de colocação pronominal ( próclise e ênclise ), distintas dos hábitos da oralidade ou da escrita informal ; d ) Seleção vocabular adequada à modalidade escrita formal da Língua Portuguesa, exigida pela situação comunicativa – essa competência envolve a precisão na utilização do vocabulário relacionado à temática solicitada pela questão ; a ausência de marcas da oralidade, como termos de sentido de muito genérico ( “ coisa ”, “ negócio ”, “ você ” ) e termos de registros mais informais ( como gírias, jargões, frases feitas, ditados populares, termos regionais ).
Assim, espera -se que o participante respeite a adequação vocabular não utilizando gírias ou expressões coloquiais, evite repetição desnecessária de palavras e utilize um vocabulário mais formal, adequado ao texto de caráter dissertativo.
A escolha dessas competências para subsidiar o processo de avaliação apoia -se na concepção de que, no desempenho dos graduandos, a modalidade escrita tem apresentado uma intensa simplificação, originada no padrão da modalidade oral da Língua Portuguesa.
No caso do texto de base dissertativa, inscrito em um registro formal, a distância entre as duas modalidades é ainda maior, o que provoca situações de hipercorreção ( desvios provocados pela incorporação indevida de uma regra da norma-padrão ) e de truncamentos sintáticos ( estruturas frasais incompreensíveis devido à complexidade sintática própria da modalidade escrita ).
Observam -se, então, os seguintes aspectos que marcam essa distinção entre as duas modalidades, devido à excessiva simplificação da modalidade falada : a ) redução drástica de estruturas subordinadas, compensada pelo aumento na frequência de estruturas coordenadas e absolutas, por um lado, ou pela elaboração de estruturas truncadas pelo excesso de ideias sem a devida conexão subordinativa ; b ) redução no uso de conectores para expressar relações lógicas essenciais à construção do texto, substituídas pela exigência de inferência por parte do interlocutor para suprir a sua ausência ; c ) redução cada vez maior do uso do subjuntivo, ao lado da ampliação do uso do indicativo combinado a estruturas frasais coordenadas ou absolutas ; d ) empobrecimento do processo de referenciação, com a repetição exaustiva de pronomes ou nomes ; e ) simplificação extrema da marcação da categoria tempo na morfologia verbal ; f ) falta de domínio de vocabulário mais abstrato e de maior complexidade, essencial ao desenvolvimento do processo dissertativo ; g ) redução drástica no emprego da acentuação gráfica, processo intensificado pela divulgação imprecisa das mudanças promovidas pelo último acordo ortográfico.
Os aspectos macroestruturais da elaboração do texto não foram avaliados neste processo, para não penalizar duplamente os estudantes, já que a banca de formação geral, composta por profissionais de diferentes áreas do conhecimento, encarregou -se da avaliação do conteúdo desenvolvido nas questões.
São eles : progressão temática, coerência na relação com os conhecimentos de mundo dos usuários da língua, inteligibilidade, atendimento ao solicitado no enunciado do ponto de vista do desenvolvimento do conteúdo, entre outros.
A grade de avaliação do desempenho linguístico considerou, portanto, três grandes grupos de competências, segundo os aspectos explicitados anteriormente : 1.
Aspectos ortográficos : domínio das convenções ortográficas : grafia de vogais e consoantes, uso de maiúsculas e minúsculas, emprego do hífen e acentuação gráfica ; 2.
Aspectos textuais : domínio dos procedimentos de estruturação textual do ponto de vista microestrutural : organização interna dos períodos, emprego de conectores para a articulação lógica entre os períodos e entre os parágrafos, emprego de marcas de referenciação lexical e pronominal ; utilização dos sinais de pontuação que contribuem para a organização lógica da frase. 3.
Aspectos morfossintáticos e vocabulares : domínio das regras de caráter morfossintático estabelecidas como modelares do ponto de vista da modalidade escrita formal da norma-padrão da Língua Portuguesa : concordância nominal e verbal, regência nominal e verbal, colocação pronominal, flexão nominal e verbal, correlação entre tempos e modos verbais, ausência de marcas de oralidade.
A seleção vocabular adequada à modalidade escrita formal da Língua Portuguesa foi incorporada a essa última competência, tendo em vista a intersecção entre as duas do ponto de vista das exigências do registro formal da modalidade escrita da norma-padrão.
Os resultados da avaliação correspondem aos seguintes aspectos observados em cada competência : Aspectos ortográficos : A correção foi realizada classificando os textos em cinco níveis, nível zero a quatro.
O desempenho dos estudantes ficou concentrado, majoritariamente, no nível 3 ( de 4 a 7 desvios ), devido, principalmente, aos desvios de acentuação.
O número de desvios de grafia foi reduzido.
Diferentemente do resultado do ENADE/2013, aumentou a porcentagem de textos sem desvios ( enquadrados no nível 4 ) e diminuiu a porcentagem de textos enquadrados no nível 2.
Alguns textos, com pior desempenho, foram enquadrados no nível 1 ( de 8 a 12 desvios ), enquanto o nível zero foi atribuído a pouquíssimos casos, já que ele revela falta absoluta de domínio das convenções ortográficas.
Observou -se, portanto, que existe uma diferença muito grande de desempenho nos dois aspectos analisados : baixo índice de desvios ortográficos e grande índice de desvios de acentuação.
Em vários casos, ocorre ausência completa de acentuação gráfica.
Os resultados revelam, portanto, que a tendência dominante entre os universitários brasileiros é a eliminação da acentuação gráfica, provavelmente motivada pela vivência dos jovens relacionada aos aplicativos de comunicação via internet ( redes sociais e e-mails ).
Nesse tipo de comunicação, devido ao ritmo intenso de troca de mensagens, o uso de acentos gráficos foi praticamente abolido.
Outro fator que pode ter relação com essa tendência é a ausência de esclarecimento dos meios de comunicação, das autoridades e das escolas sobre as decisões do Acordo Ortográfico da Língua Portuguesa de 1990, gerando um estado de indefinição para os estudantes.
Os casos mais sistemáticos de eliminação do acento indicador da sílaba tônica são :  palavras proparoxítonas ( “ ridiculos ”, “ publicas ”, “ lideres ”, “ politicos ” ) ;  palavras paroxítonas terminadas em ditongo crescente ( “ homicidios ”, “ latrocinio ”, “ individuo ”, “ dependencia ” ) ;  palavras oxítonas ( “ ninguem ”, “ esta ”, “ ate ”, “ ai ” ).
Por outro lado, destaca -se o uso indevido do acento gráfico em determinadas palavras, como observado nas grafias * “ jornáis ”, * “ telejornaís ”, * “ propíciar ”, * “ medídas ”, * “ dígnidade ”, * “ cídades ”.
Quanto ao domínio das convenções relativas à grafia das palavras, observam -se desvios como : a hipercorreção pela escolha de “ e ” no lugar de “ i ”, por influência de hábitos da oralidade ( * “ enumeros ” por “ inúmeros ”, * “ entevenção ” por “ intervenção ” ) ; a eliminação do “ r ” marcador do infinitivo verbal ( * “ esta ” no lugar de “ estar ” ).
Outros casos de desvios de grafia relacionados à variação diastrática podem ser observados em *“estrupo ”, *“automovís ”, *“viensse ”, *“camihada ”.
São muito frequentes os seguintes desvios de caráter ortográfico, com repercussão morfossintática :  eliminação da marca de infinitivo ( -r- ) e substituição por acento agudo ( “ está ” no lugar de “ estar ” ) ;  confusão entre “ ão ” e “ am ” nas formas verbais ( “ invadão ” no lugar de “ invadam ” e “ estam ” no lugar de “ estão ” no presente do indicativo ; “ estaram ” no lugar de “ estarão ” no futuro do indicativo ) ;  confusão entre a grafia do verbo “ haver ” ( “ há ” ) e o artigo definido ou a preposição “ a ” ;  uso de hífen para separar pronome átono – tanto uso indevido quanto omissão ( no pretérito imperfeito e futuro do subjuntivo : “ evitar -mos ” no lugar de “ evitarmos ”, “ percebesse ” no lugar de “ percebe -se ” e vice-versa ) ; Observam -se, também, muitos casos de inadequação no uso da maiúscula : ausência de diferença entre a primeira letra e as outras, em início de período, principalmente dos estudantes que adotam a escrita em letra de imprensa ; utilização de maiúscula para destacar determinadas palavras-chave do texto, como “ Violência ”, “ Brasileiros ”, “ Fatores ”, “ Ozônio ”, “ Sustentável ”.
Vale observar, também, que, ao contrário do que se esperava, não apareceram abreviaturas próprias do “ internetês ”, ou seja, dos hábitos de comunicação escrita adquiridos pelo uso de redes sociais e emails.
Aspectos textuais : Quanto a esses aspectos a correção também classificou os textos em cinco níveis ( zero a quatro ), em função da quantidade de erros apresentados.
O desempenho dos estudantes ficou concentrado, majoritariamente, nos níveis 3 e 2, devido à grande ocorrência de problemas de estruturação textual.
Foi muito baixo o número de textos que não apresentaram qualquer problema estrutural e, portanto, ficaram enquadrados no nível 4.
Alguns textos com pior desempenho foram enquadrados no nível 1, enquanto o nível zero foi atribuído a pouquíssimos casos, já que ele revela a existência de um texto sem articulação e com comprometimento do sentido.
Observou -se que a grande maioria dos estudantes não distribuiu as ideias em parágrafos, talvez devido ao pequeno número de linhas disponibilizadas para a resposta da questão ou, quem sabe, pela suposição de que não seria necessária essa divisão por não se tratar de um texto no modelo de uma redação dissertativo-argumentativa, como solicitado nos vestibulares.
Outro aspecto observado na estruturação textual foi a divisão em dois itens, provavelmente motivados pelo encaminhamento do enunciado das questões, que apresentavam dois pontos a serem detalhados : a questão 1 solicitava que os estudantes dissertassem sobre as consequências do transporte motorizado ( a ) e apresentassem ações de intervenção por parte do poder público ( b ) ; a questão 2 solicitava que os estudantes dissertassem sobre as causas da violência urbana ( a ) e os fatores para evitá -la ( b ).
Essa competência pode ser considerada como a mais problemática no que diz respeito ao desempenho linguístico dos estudantes, porque são muitos os problemas observados, desvios acumulados durante toda a formação escolar e que não se resolvem com um estudo autodidata, como acontece com regras ortográficas ou morfossintáticas : sequência justaposta de ideias sem encaixamentos sintáticos ; redução drástica de estruturas subordinadas, ao lado do aumento na frequência de estruturas coordenadas e absolutas ; redução no uso de conectores para expressar relações lógicas essenciais à construção do texto, substituídas pela exigência de inferência por parte do interlocutor para suprir a sua ausência ; emprego equivocado de operadores que não estabelecem relações lógicas coerentes entre ideias do texto ; emprego inadequado do pronome relativo ( com omissão da preposição ou a utilização de pronome inadequado, como “ onde ” ) ; repetição exaustiva de termos sem a utilização de procedimentos mais sofisticados de substituição ( hiperonímias, hiponímias, nominalizações, expressões metafóricas ) ; frases fragmentadas que comprometem a estrutura lógico-gramatical ; frases formadas apenas por oração subordinada, sem oração principal.
Um importante aspecto a destacar é o baixíssimo desempenho de uma parte dos estudantes em relação à estrutura formal do texto produzido, o que é extremamente preocupante ao se levar em conta que são graduandos em fase final de formação.
São frequentes os casos de desvios de estruturação frasal, com uso inadequado ou ausência de conectivos entre parágrafos e entre frases.
Em uma parte dos textos, falta um mínimo de textualidade e de domínio do registro padrão da língua.
Na verdade, observam -se relações linguísticas quase agramaticais, como as estabelecidas pela sequência de gerúndios sem o apoio de um ponto de partida para a organização das informações gramaticais e semânticas, ou seja, sem uma oração principal.
Quanto à utilização dos mecanismos de referenciação, deve -se destacar a ocorrência, em uma boa parte dos textos, de repetições de palavras ou expressões sem a utilização de termos sinônimos ou pronomes, como seria adequado.
Quanto à utilização dos sinais de pontuação, observou -se uma grande precariedade nos textos analisados.
É muito frequente a ocorrência de parágrafos sem marca interna de pontuação para separar os períodos.
Vale observar que não foi penalizada a ausência de vírgula para destacar locuções ou adjuntos adverbiais de pequena extensão deslocados de posição na frase, por ser um uso opcional.
São os seguintes os tipos de problemas encontrados :  vírgula : utilização de vírgula para separar o sujeito e o predicado ; ocorrência de apenas uma das vírgulas para separar uma palavra, uma expressão ou uma oração encaixada ; uso de vírgula no lugar do ponto para separar ideias que constituem períodos distintos ; ausência de vírgula para separar elementos de uma enumeração ; ausência de vírgula para  separar oração adjetiva explicativa ou utilização inadequada para separar oração adjetiva restritiva ;  ponto e vírgula : utilização do ponto e vírgula no lugar de vírgula ;  ponto final : ausência de ponto final para separar períodos.
Aspectos morfossintáticos e vocabulares : Da mesma forma que nos aspectos anteriores os textos foram classificados em níveis de zero a quatro.
O desempenho dos estudantes ficou concentrado, majoritariamente, nos níveis 3 e 2, devido à grande ocorrência de problemas de regência e concordância.
O nível 4 foi atribuído a um número menor de textos.
Alguns textos com pior desempenho foram enquadrados no nível 1, enquanto o nível zero foi atribuído a pouquíssimos casos, já que ele revela a existência de um texto sem o respeito às mínimas exigências morfossintáticas da norma-padrão e com comprometimento do sentido.
Os resultados são muito transparentes em relação aos aspectos mais problemáticos do desempenho dos estudantes.
O desvio mais frequente, em relação à regência, é a falta do sinal indicativo da crase – isso revela que o usuário não tem consciência de que, sob a forma do termo “ a ”, existe a presença de uma preposição “ a ”, exigida pela regência do termo anterior.
Embora em outros exames, como o Enem, a falta de crase seja penalizada em acentuação, nesta correção esse desvio foi considerado no âmbito dos aspectos morfossintáticos.
Outro problema relacionado à regência verbal e à nominal, encontrado frequentemente nas questões, foi a ausência de preposição antes de pronome relativo, processo generalizado na modalidade oral da língua, em situações de registro informal.
Apesar da possibilidade de que essa alteração de regência se generalize no padrão escrito da Língua Portuguesa, como já está ocorrendo até em textos jornalísticos, o não emprego da preposição foi penalizado neste processo de avaliação.
Outro desvio muito frequente diz respeito aos processos de concordância verbal e de concordância nominal.
Quanto à concordância de número, observou -se ausência de marca ( com sujeito anteposto ou posposto ) ou uso indevido ( uso inadequado da marca de plural comandado pelo núcleo plural da locução adjetiva, apesar de o substantivo que funciona como núcleo do sintagma nominal estar no singular ).
Uma ocorrência que se destacou foi a ausência de acento circunflexo na forma plural do presente do indicativo dos verbos “ ter ” e “ vir ”, que foi considerada como um desvio na concordância verbal e não na acentuação gráfica.
Quanto à concordância de gênero, vários casos foram observados, normalmente no âmbito de sintagmas nominais longos, em que o adjetivo está afastado do substantivo.
Deve -se destacar uma ocorrência não observada no ENADE/2013 : o aparecimento da marca de plural em verbos ou adjetivos relacionados a núcleos substantivos no singular, evidenciando um processo de hipercorreção.
Quanto à questão da colocação pronominal, foram poucos os casos observados.
Apesar de serem aspectos relacionados à oralidade, concluiu -se que, no registro escrito formal, a maioria dos estudantes já incorporou regras como a não introdução da frase por um pronome oblíquo e a próclise na presença de um termo atrator.
Não se adotou, entretanto, o padrão excessivamente formal descrito pelas gramáticas normativas em relação à posição do pronome oblíquo em locuções verbais, já que esse uso está muito distante da prática cotidiana, até em textos mais formais.
Quanto aos aspectos vocabulares, alguns tipos de inadequação foram observados : expressões da oralidade apareceram em algumas respostas, mas sem maior relevância do ponto de vista quantitativo ; seleção vocabular incompatível com o contexto, gerando situações de falta de inteligibilidade ; falta de domínio de vocabulário mais abstrato e de maior complexidade, essencial ao desenvolvimento do texto de base dissertativa.
Várias marcas de oralidade foram identificadas, embora não em alta frequência : o uso do pronome relativo “ onde ” como relativo universal, falta de artigo definido antes de substantivo, repetição de palavras por falta de vocabulário, reduções como “ tá ”, “ pra ”, “ pro ”, “ prum ”, expressões informais. 3.3.2 Componente de Conhecimento Específico Na parte da prova relativa às questões discursivas no Componente de Conhecimento Específico ( Tabela 3.17 ), observa -se que a média foi bem mais baixa do que para as questões discursivas do Componente de Formação Geral.
Enquanto no Componente de Formação Geral a média para estudantes de Letras - Português ( Bacharelado e Licenciatura ) de todo o Brasil foi 52,6, na parte de Conhecimento Específico a média foi 25,3.
A maior média deste componente foi obtida pelos estudantes da região Sudeste ( 27,5 ), e a menor, pelos da região Centro-Oeste ( 23,3 ).
Quanto à variabilidade das notas, o desvio padrão de todo o Brasil foi 16,6.
O maior desvio padrão foi encontrado na região Sudeste ( 18,7 ), e o menor, nas regiões Nordeste e Sul ( 15,4 ). concordância de gênero, vários casos foram observados, normalmente no âmbito de sintagmas nominais longos, em que o adjetivo está afastado do substantivo.
Deve -se destacar uma ocorrência não observada no ENADE/2013 : o aparecimento da marca de plural em verbos ou adjetivos relacionados a núcleos substantivos no singular, evidenciando um processo de hipercorreção.
Quanto à questão da colocação pronominal, foram poucos os casos observados.
Apesar de serem aspectos relacionados à oralidade, concluiu -se que, no registro escrito formal, a maioria dos estudantes já incorporou regras como a não introdução da frase por um pronome oblíquo e a próclise na presença de um termo atrator.
Não se adotou, entretanto, o padrão excessivamente formal descrito pelas gramáticas normativas em relação à posição do pronome oblíquo em locuções verbais, já que esse uso está muito distante da prática cotidiana, até em textos mais formais.
Quanto aos aspectos vocabulares, alguns tipos de inadequação foram observados : expressões da oralidade apareceram em algumas respostas, mas sem maior relevância do ponto de vista quantitativo ; seleção vocabular incompatível com o contexto, gerando situações de falta de inteligibilidade ; falta de domínio de vocabulário mais abstrato e de maior complexidade, essencial ao desenvolvimento do texto de base dissertativa.
Várias marcas de oralidade foram identificadas, embora não em alta frequência : o uso do pronome relativo “ onde ” como relativo universal, falta de artigo definido antes de substantivo, repetição de palavras por falta de vocabulário, reduções como “ tá ”, “ pra ”, “ pro ”, “ prum ”, expressões informais. 3.3.2 Componente de Conhecimento Específico Na parte da prova relativa às questões discursivas no Componente de Conhecimento Específico ( Tabela 3.17 ), observa -se que a média foi bem mais baixa do que para as questões discursivas do Componente de Formação Geral.
Enquanto no Componente de Formação Geral a média para estudantes de Letras - Português ( Bacharelado e Licenciatura ) de todo o Brasil foi 52,6, na parte de Conhecimento Específico a média foi 25,3.
A maior média deste componente foi obtida pelos estudantes da região Sudeste ( 27,5 ), e a menor, pelos da região Centro-Oeste ( 23,3 ).
Quanto à variabilidade das notas, o desvio padrão de todo o Brasil foi 16,6.
O maior desvio padrão foi encontrado na região Sudeste ( 18,7 ), e o menor, nas regiões Nordeste e Sul ( 15,4 ).
A maior nota máxima foi obtida na região Sudeste ( 90,0 ), enquanto a menor nota máxima foi encontrada na região Centro-Oeste ( 75,0 ).
Além disso, a nota mínima ( 0,0 ) foi obtida por pelo menos um estudante em todas as regiões do Brasil, sem exceção.
A mediana do Brasil como um todo foi 23,3.
Nas regiões Norte, Nordeste e Centro-Oeste a também mediana foi 23,3, e nas regiões Sudeste e Sul, a mediana foi 26,7 e 25,0 respectivamente.
O Gráfico 3.11 representa a distribuição das notas nas questões discursivas no Componente de Conhecimento Específico segundo a opção de Bacharelado/Licenciatura.
As modas destas distribuições ocorrem no primeiro intervalo, [ 0 ; 10 ] para o Bacharelado, e no intervalo ( 20 ; 30 ] para a Licenciatura.
A análise de cada uma destas questões será feita a seguir. 3.3.2.1 Análise da Questão Discursiva 3 do Componente de Conhecimento Específico Para a questão 3, presente nas provas de Bacharelado e Licenciatura, os resultados aferidos encontram -se descritos na Tabela 3.18.
O desempenho dos estudantes de todo o Brasil nesta questão foi o mais baixo dentre as três questões discursivas desse componente, e a média dos estudantes de todo o Brasil foi 12,3.
A menor média nessa questão foi obtida pelos estudantes das regiões Norte e Nordeste ( 10,5 em ambas ), enquanto a maior média foi obtida na região Sudeste ( 16,9 ).
Quanto à variabilidade das notas, o desvio padrão de todo o Brasil foi 19,7.
O maior desvio padrão foi obtido na região Sudeste ( 23,6 ), enquanto o menor foi obtido na região Nordeste ( 17,4 ).
A nota máxima, 100,0 pontos, foi alcançada por, pelo menos, um estudante em todas as regiões.
A mediana e a notas mínima foram zero em todas as regiões do Brasil, sem exceção, o que significa que pelo menos 50 \% dos estudantes obteve nota zero ou deixou a questão em branco.
O Gráfico 3.12 mostra a distribuição das notas na questão discursiva 3 do Componente de Conhecimento Específico da área de Letras - Português para estudantes do Bacharelado e da Licenciatura.
Essa distribuição tem moda nas questões em branco, com percentual maior para os estudantes de Licenciatura. 3.3.2.2 Comentários sobre as respostas à Questão Discursiva 3 A questão, ao tratar do conceito do Belo na arte produzida na Antiguidade Clássica e cotejá -lo com manifestações tanto arquitetônicas quanto literárias da Idade Média, satisfaz inteiramente os incisos I e IV do Art. 5º da Portaria Inep nº 258, de 02/06/2014, publicada no D.O.U. de 04/06/2014 : “ I - Competência intercultural ; [... ] e IV- capacidade de analisar criticamente diferentes manifestações linguísticas e literárias ”.
Essa questão, igualmente, preenche a habilidade prevista no Art. 6º, inciso IV “ refletir analítica, crítica e teoricamente acerca das manifestações literárias e culturais ” bem como abarca o conteúdo especificado no Art. 7º, inciso II.a “ conceitos de literatura e arte ”.
Tal adequação também estava presente no enunciado que não deixava margem a dúvidas sobre os conceitos pedidos nem sobre a habilidade que os estudantes necessitavam ter para responder à questão, chegando mesmo a facilitar a resposta ao ilustrar a questão com a imagem de uma catedral gótica.
O padrão de respostas apresentava o que se espera de um estudante mediano de Letras, ou seja, exigia o básico do conhecimento previsto sobre a Antiguidade Clássica e a Idade Média, nada além do que costuma estar previsto nos currículos de graduação.
Por isso pode -se classificar a questão como fácil.
De fato, houve respostas que alcançaram plenamente o padrão e outras que até mesmo chegaram a inserir mais informações, como aquelas que relacionaram o padrão da época Clássica a obras do Renascimento, com menção a artistas e obras ( Da Vinci, a Pietà, entre outras ).
Por outro lado, houve também estudantes que mal conseguiram parafrasear o enunciado e não associaram a verticalidade da arquitetura gótica ao texto, como expressão de ascensão da alma.
Os equívocos cometidos, como o de desconhecer que a noção de Belo não é absoluta, e as tentativas malsucedidas de paráfrase, são frutos tão somente da não familiaridade do estudante com assunto objeto da questão.
Chamou a atenção, sobretudo, o abandono do texto de “ A Demanda do Santo Graal ”.
Registrou -se quantidade expressiva de estudantes que abordou ( ou que, pelo menos, tentou abordar ) a noção de Belo no Classicismo e na Idade Média, mas sequer tangenciou o assunto da morte da Galaaz.
Pareceu, inclusive, que alguns estudantes não se deram conta de que o texto narrava uma morte, ou seja, não perceberam o caráter alegórico do trecho “ a alma lhe saiu do corpo e levavam -na os anjos ”.
No que concerne à linguagem utilizada, essa variou desde o emprego demasiadamente precário do vernáculo ao domínio pleno da norma padrão da língua.
A média geral das notas atribuídas a essa questão foi baixo, classificado como fraco para ambas as habilitações, sendo que os do Bacharelado obtiveram uma média mais alta e próxima ao dobro do valor da média obtida pelos de Licenciatura.
A questão 3 registrou um percentual elevado de respostas em branco, evidenciando que uma quantidade expressiva dos estudantes não se sentiu motivado a tentar redigir uma resposta, possivelmente por exigir a associação de conhecimentos interdisciplinares do campo da Literatura, da Arte e da História.
Evidentemente, a diversidade de erros foi encontrada com maior frequência nas respostas medianas e fracas, já que as respostas boas apresentaram um menor grau de variação de abordagem.
A questão foi composta de três itens, mas independentemente do fato de esses itens terem sido redigidos na forma de um único texto ou de virem itemizados, os estudantes produziram respostas fracas ou, no máximo, medianas.
Se dominava apenas um dos itens, este era normalmente o concernente ao Belo no Classicismo.
Os conceitos e valores que imperavam na Idade Média – o teocentrismo e, sobretudo o gótico – estavam expressos de modo muito confuso ou estavam erradamente descritos.
Aliás, vale notar que há diferença entre o que os estudantes revelavam saber sobre cada um desses tópicos.
Por exemplo, a oposição teocentrismo x antropocentrismo aparecia com frequência considerável, e, muitas vezes, era o que evitava o zero.
Já a associação entre Belo e simetria/proporções áureas/corpo humano, ocorreu muito raramente.
O mais difícil para os estudantes, como já salientado, foi estabelecer a relação entre o texto e a catedral, o que exigia um esforço maior de abstração.
A correção das respostas revelou que há diferença qualitativa entre o desempenho dos estudantes do curso de Bacharelado e o dos estudantes de Licenciatura.
De um modo geral, as respostas dos primeiros revelam maior consistência de formação e domínio dos conteúdos, competências e habilidades descritas nas Diretrizes Curriculares.
Tal diferença fica nítida não apenas na correção da questão 3 ( Literatura ), mas também na da questão 4 ( Análise Linguística ), questões comuns às duas habilitações.
Os estudantes de Bacharelado foram muito mais precisos e críticos em ambas as questões do que os de Licenciatura.
Essa diferença se registrou até mesmo na qualidade do erro.
Em relação à questão 3, de um modo geral, os estudantes de Bacharelado erram revelando alguma noção, mesmo que longínqua, do que está sendo pedido ; os de Licenciatura erram demonstrando desconhecimento do assunto.
Destaca -se que o erro mais frequentemente observado foi o de interpretação do texto e a relação pouco observada do texto com a arquitetura gótica.
Atribui -se tal fato ao nível de abstração requerido, o qual boa parte dos estudantes pareciam não estar habilitados a alcançar. 3.3.2.3 Análise da Questão Discursiva 4 do Componente de Conhecimento Específico A Tabela 3.19 contém as informações relativas à questão 4 do conjunto de questões do Componente de Conhecimento Específico dos concluintes de Bacharelado e Licenciatura em Letras-Português.
A média geral do Brasil foi 23,4, sendo a menor média registrada na região Centro-Oeste ( 21,9 ) e a maior, na região Sudeste ( 25,3 ).
As medianas ( 20,0 ), as notas máximas ( 100,0 ) e as notas mínimas ( 0,0 ) foram as mesmas para todas as regiões, exceto a nota máxima na região Centro-Oeste ( 90,0 ).
O Gráfico 3.13 representa a distribuição de notas na questão discursiva 4, no Componente de Conhecimento Específico para estudantes do Bacharelado e da Licenciatura.
A maior frequência, tanto para estudantes de Bacharelado quanto para Licenciatura, ocorre na categoria “ em branco ”. 3.3.2.4 Comentários sobre as respostas à Questão Discursiva 4 A questão, ao tratar da entrada do referente limão, e sua consequente lexicalização, na língua indígena Guató, cobrou do estudante conhecimentos como contato linguístico e mudança, trocas culturais e sua interveniência na competência e desempenho de falante nativo, além de causas de desaparecimento de uma dada língua, incluídos níveis de análise gramatical ( morfologia e fonologia ).
Assim, a questão estava associada aos incisos II e IV do Art. 5º da Portaria Inep ( op. cit. ), que prevê o perfil do profissional da área.
No que concerne às competências e habilidades necessárias, a questão se enquadrava perfeitamente no Art. 6º, inciso II ; no que diz respeito aos conteúdos curriculares, abordava consistentemente o determinado pelo Art. 7º, incisos I.a e I.b.
Dessa forma, a questão era absolutamente compatível com o que se espera de um profissional da Área de Letras.
A questão exigia, de fato, conhecimentos técnicos de fonologia e de morfologia, porque, embora o texto 2 explicasse e descrevesse fatos morfofonológicos do Guató, Essa descrição era feita em linguagem e notações técnicas, o que requeria do estudante capacidade de entendimento da regra.
Isso não significa que faltou clareza, ao contrário, a linguagem técnica não tem ambiguidade.
Mas, foram poucos os que realmente entenderam a questão.
Quando tentavam responder, via de regra, faziam uma paráfrase ( quando não uma cópia ) do enunciado.
Muitas vezes, até mesmo a paráfrase estava errada.
Para responder ao item ‘ b ’, parte significativa dos estudantes atribuiu o desaparecimento das línguas à inexistência de pesquisa linguística, como se as operações de registro, descrição e análise pudessem manter a língua viva.
O que tais operações podem fazer é manter vivo o conhecimento da língua, não a sua produtividade.
Esse ponto provocou digressões e equívocos nas respostas, comprometendo a qualidade da maioria delas.
O nível de dificuldade da questão é função da qualidade do curso que os estudantes fizeram e do modo como o fizeram.
Pressupondo um estudante mediano, que frequentou um curso em que tenha havido formação em morfologia e fonologia, a questão é de média dificuldade para fácil ; caso contrário, dificílima, principalmente no tocante ao item ‘ a ’.
O padrão de respostas não exigia considerações técnicas sobre os fenômenos morfofonológicos, o que facilitou não só a correção, como também elevou o nível das notas.
Quando os estudantes chegavam a mencionar especificidades técnicas ( e houve aqueles que o fizeram ), via -se a qualidade de sua formação.
A análise das respostas dadas ao item ‘ a ’ permite agrupá -las segundo algumas ocorrências gerais : • O espaço reservado a esse item foi deixado em branco ; • O estudante fez uma paráfrase do enunciado no que diz respeito à causa da adoção da palavra pelo Guató ; • O estudante conseguiu discutir a possibilidade de compatibilização entre os dois padrões morfofonológicos ( o do Guató e o do Português ).
Em relação ao item ‘ b ’, as respostas podem assim ser agrupadas : • O estudante atribuiu o desaparecimento das línguas à extinção dos povos indígenas ; • O estudante atribuiu desaparecimento das línguas a doenças contraídas pelos índios em decorrência do contato com os “ civilizados ”.  O estudante responsabilizou os pesquisadores pela falta de interesse da linguística pelas línguas em extinção numa clara inversão da questão – a língua desapareceria simplesmente porque não é estudada.  O estudante atribuiu o desaparecimento das línguas ao deslocamento geográfico da população, o que impacta a cultura, ocasionando a perda da língua original.
Registre -se, entretanto, que houve respostas completas, atendendo ao padrão esperado de respostas.
O enunciado dessa questão era claro, não ensejando ambiguidades de leitura.
Os erros cometidos em relação ao item ‘ a ’, dado seu caráter técnico, deveram -se à falta de conhecimento do assunto pelo estudante.
As paráfrases, dependendo do nível de aproximação do problema, podem demonstrar algum conhecimento sobre o tema ou podem, tão somente, ser meras tentativas de registrar algum conteúdo que pudesse ser convertido em acréscimo de pontuação.
Diferentemente do item ‘ a ’, o item ‘ b ’ oferecia, por \textbf{intermédio} do texto 2, boas “ pistas ” para ativar o conhecimento do estudante sobre a causa de desaparecimento das línguas indígenas.
Por outro lado, mencionava a falta de pesquisa, o que, por sinal, não está presente no padrão de respostas.
O estudante que desconhece questões de contato, mudança, pressão de cultura dominante, etc., toma esse fato, que não está relacionado com a extinção das línguas, como sua causa, sem refletir que a consequência da ausência de pesquisa é, no máximo, a língua ficar desconhecida por falta de descrição.
Dados os fatores acima comentados, o desempenho geral dos estudantes pode ser considerado fraco, como mostram as médias das notas atribuídas às respostas dadas às questões.
Nas duas habilitações o percentual de resposta em branco foi em torno de 30 \%.
Houve expressiva diversidade nas soluções apresentadas, algumas muito inesperadas como : a proposição de leis que obrigassem ao ensino das línguas indígenas ; o comentário de que o contato linguístico trouxe como consequência a miscigenação das raças ( invertendo a questão ) ; a adoção da abordagem histórica seguida de observação de influências lexicais, entre outras.
Foram frequentes as respostas genéricas, que demonstravam falta de conhecimento técnico da questão, baseadas exclusivamente no senso comum, como as que afirmavam que o desaparecimento dos índios ocorreria devido à falta de ajuda das pessoas ou à ausência de políticas públicas.
A principal abordagem foi a do tipo genérico, sem base conceitual, e o principal erro registrado consistiu em atribuir a responsabilidade do desaparecimento das línguas à falta de pesquisa.
Registre -se, também, como fato preocupante, a incapacidade demonstrada por boa parte dos estudantes de lidar com as regras morfofonológicas.
Sem dúvida, pode -se afirmar que os componentes Fonologia e Morfologia não estão sendo suficientemente estudados nas universidades de um modo geral. 3.3.2.5 Análise da Questão Discursiva 5 do Componente de Conhecimento Específico ( Bacharelado ) A Tabela 3.20 contém as informações relativas à questão 5 do conjunto do Componente de Conhecimento Específico ( Bacharelado ).
O desempenho dos estudantes nessa questão foi um pouco superior ao das questões 3 e 4.
A nota média dos estudantes de todo o Brasil foi 29,6.
A maior média foi registrada na região Sudeste ( 32,1 ), enquanto a menor média foi registrada na região Norte ( 22,7 ).
Quanto à variabilidade das notas, o desvio padrão dos estudantes do Brasil, como um todo, foi 29,8.
Enquanto o maior desvio foi encontrado na região Centro-Oeste ( 30,0 ), o menor foi encontrado na região Norte ( 27,1 ).
A mediana das notas dos concluintes de todo o Brasil foi 30,0.
A menor mediana foi zero, a das regiões Norte e Centro-Oeste.
A maior mediana ocorreu na região Sudeste ( 30,0 ).
A nota máxima ( 100,0 ) foi obtida por pelo menos um estudante das regiões Nordeste, Sudeste e Centro-Oeste.
Já na região Norte a nota máxima foi 55,0 e na região Sul foi 90,0.
A nota mínima foi zero para todas as regiões, sem exceção.
O Gráfico 3.14 representa a distribuição das notas da questão discursiva 5 no Componente de Conhecimento Específico ( Bacharelado ).
Mais uma vez, destaca -se o grande número de estudantes que deixaram a questão em branco, categoria correspondendo à moda. 3.3.2.6 Comentários sobre as respostas à Questão Discursiva 5 ( Bacharelado ) A questão a ser respondida pelos estudantes oriundos de cursos de Bacharelado, avaliava os conteúdos previstos nos incisos II.a ; II.b e II.h do Art. 7º da Portaria Inep e as competências e habilidades especificadas nos incisos II e IV do Art. 6º do mesmo documento.
A questão se estruturava a partir de dois textos que facilitavam a resposta dos estudantes.
O trecho do texto de Mia Couto, autor bem divulgado no Brasil, além de marcas linguísticas que caracterizam a oralidade do Português moçambicano ( como “ controlo ”, de forte influência do Português europeu ) apresenta fenômenos que poderiam contribuir para o estudante estabelecer relação com as manifestações literárias de outras variedades do Português, como neologismos, aliterações, caracterização de personagem, etc.
Sendo assim, em sua resposta, o estudante poderia buscar aproximações com autores brasileiros, como Guimarães Rosas e, principalmente, Graciliano Ramos, entre outros.
O texto teórico, adaptado, oferecia dados para que o estudante pudesse estabelecer a relação solicitada sobre literatura e cultura.
A boa seleção dos textos que antecederam o comando evocou a lembrança de autores brasileiros, facilitando ao estudante citá -los com segurança.
O comando era claro e a questão não oferecia qualquer ambiguidade.
Ao contrário, a boa formulação do enunciado permitiu ao estudante expressar seu conhecimento sobre o tema focalizado e demonstrar amadurecimento quanto aos pontos selecionados.
Houve boa compatibilidade entre o padrão de resposta e as respostas dadas pelos estudantes, principalmente no que diz respeito à relação entre língua e cultura, literatura como representante de cultura.
Algumas respostas destacaram também a relação da literatura com os usos linguísticos.
Embora conciso em relação a sua apresentação, o caráter abrangente do padrão de respostas contemplou um amplo conjunto de possibilidades de solução.
De todas as questões corrigidas, essa, especificamente dirigida aos Bacharelandos, foi aquela cujas respostas foram mais bem pontuadas, registrando a melhor média dentre as médias do Bacharelado.
No entanto, dentre os estudantes de bacharelado, cerca de 10 \% receberam nota zero para a resposta apresentada e pouco mais do que 30 \% deles não responderam a essa questão.
Foram encontradas respostas fracas como as que apenas mencionaram superficialmente a existência da relação entre literatura e cultura ou que descreveram a literatura como reflexo de uma sociedade ( ou ainda outras “ criativas ” como a que atribui à personagem principal o fato de ser indiana por causa das relações comercias Índia/Moçambique ).
Foram encontradas, também, respostas que demonstravam bastante proficiência sobre a literatura africana e que identificaram na obra de Mia Couto a presença de elementos telúricos, relacionando esses elementos a obras de outros autores africanos, como Agualusa.
Vale notar que também foram encontradas observações sobre as literaturas lusófonas que, embora utilizem a mesma língua, apresentam realizações diferentes por causa de suas variantes e diferentes representações da realidade, decorrentes do modo de organização das sociedades que se encontram na obra retratadas.
O aspecto mais bem explanado nas respostas a essa questão foi o da relação entre literatura e cultura.
A questão da presença da oralidade, que ocorre no texto de Mia Couto, bem como a presença de neologismo, proporcionaram aos estudantes a oportunidade de invocar em suas respostas referências a Guimarães Rosa, Mário de Andrade, Oswald de Andrade, etc.
Associadas também a esses aspectos, apareceram nas respostas referências à literatura brasileira nordestina, como a de Graciliano Ramos, que, de alguma forma, representa fatores histórico-culturais da literatura brasileira, bem como José Lins do Rego.
Da análise comparativa dos resultados alcançados pelos estudantes em diferentes questões, pode -se deduzir que a área de Literatura é mais bem estudada do que a área de Língua Portuguesa.
Cabe alertar para o fato de que o texto produzido pelos estudantes, tanto em resposta a essa questão como em resposta às demais questões do componente específico, revelaram a existência de um quadro preocupante relacionado ao uso da língua, que se reflete na dificuldade de exposição do pensamento.
No caso específico dessa questão, muitos estudantes demonstraram domínio de conceitos específicos da área, como “ construção de personagem ”, “ elementos caracterizadores de paisagem ”, etc. 3.3.2.7 Análise da Questão Discursiva 5 do Componente de Conhecimento Específico ( Licenciatura ) A Tabela 3.21 contém as informações relativas à questão 5 do conjunto do Componente de Conhecimento Específico ( Licenciatura ).
O desempenho dos estudantes nessa questão foi superior ao das questões 3 e 4.
A nota média dos estudantes de todo o Brasil foi 40,8.
A maior média foi registrada na região Sul ( 42,9 ), enquanto a menor média foi registrada na região Norte ( 38,2 ).
Quanto à variabilidade das notas, o desvio padrão dos estudantes do Brasil, como um todo, foi 23,3.
O maior desvio foi encontrado na região Sudeste ( 24,5 ) e o menor foi encontrado na região Sul ( 21,7 ).
A mediana das notas dos concluintes de todo o Brasil foi 40,0, a mesma de todas as regiões.
As notas máximas ( 100,0 ) e as notas mínimas ( 0,0 ) foram iguais em todas as regiões, sem exceção.
O Gráfico 3.15 apresenta a distribuição das notas da questão discursiva 5 no Componente de Conhecimento Específico ( Licenciatura ).
Observa -se que o número de estudantes de Licenciatura que deixaram a questão 5 em branco foi o menor dentre as três questões discursivas de Conhecimento Específico.
A moda corresponde ao intervalo ( 40 ; 50 ] e a frequência de notas no intervalo ( 30 ; 40 ] é bem próxima da modal. 3.3.2.8 Comentários sobre as respostas à Questão Discursiva 5 ( Licenciatura ) Essa questão, da prova de Licenciatura, no que concerne ao perfil do profissional de Letras ( Art. 5º - Portaria Inep, op. cit. ) atendia aos incisos V, VI, VII e VIII.
Quanto às competências e habilidades requeridas, conforme especificado no Art. 6º do mesmo documento, a questão satisfazia aos incisos VI, VII, IX e X ; e em referência ao Art. 7º, os conteúdos abordados estão previstos nos incisos III.a e b.
A questão estava fundamentada em dois textos de apoio : o primeiro trazia a ilustração de um conjunto de pássaros, do qual se destaca um filhote que pia em linguagem computacional ( @pio ).
O segundo texto, teórico, antevê o impacto das tecnologias informáticas sobre o aprendizado e o desenvolvimento cognitivo da criança.
O comando da questão era claro e não permitia ambiguidade de interpretação.
O padrão de respostas contemplava a variedade de respostas dadas à questão ; e o critério de pontuação permitiu o aproveitamento máximo dos conteúdos produzidos pelos estudantes.
Para o item ‘ a ’, o padrão de respostas previa a contextualização do uso das tecnologias na sociedade, relacionando -o à formação do futuro professor, esperando que fossem mencionados variados tipos de recursos tecnológicos.
Em relação ao item ‘ b ’, esperava -se que o estudante mencionasse, pelo menos, “ dois efeitos da contínua formação profissional do professor, tais como : mudança de concepção de sala de aula e de ensino de língua, aprimoramento de recursos metodológicos para a sala de aula, facilitação do diálogo professor-aluno, socialização de práticas interescolares, promoção da autonomia do aluno, entre outros efeitos pertinentes ”.
O modelo de resposta mais frequentemente encontrada foi o que destacava a necessidade de formação continuada pelo professor.
Em relação ao item ‘ a ’, verificou -se que boa parte dos estudantes apresentou dificuldades em contextualizar “ o uso de tecnologias na sociedade, relacionando -o à formação do futuro professor de língua ”, como requerido pelo padrão de respostas.
Alguns estudantes, contudo, produziram respostas bastante aderentes ao estabelecido no padrão, evocando o uso da informática em diferentes suportes, ( como em celulares e tablets ) e em suas múltiplas aplicações ( como em jogos educativos e pesquisa escolar ).
A aplicação de tais recursos à formação do professor foi mais raramente observada.
Via de regra, os estudantes declaravam que os professores estavam mais defasados que seus alunos no uso dos recursos tecnológicos. no uso dos recursos tecnológicos.
Reconheciam que o professor envolvido em formação continuada seria mais eficiente na aplicação plena dessa tecnologia em benefício do estudante, mas mencionavam como obstáculo o desinteresse político e a falta de recursos públicos para a educação.
Em relação ao item ‘ b ’, a correção das respostas indicou que os efeitos da contínua formação do professor na prática escolar mencionados pelos estudantes se distanciavam bastante do fixado pelo padrão de resposta.
Com efeito, mostravam uma preocupante aproximação com o senso comum.
A formação contínua do professor geraria como consequência “ alunos mais bem preparados ” e proporcionaria “ melhor diálogo entre professor e aluno ”.
Sobre a promoção da autonomia do aluno, um aspecto central do debate contemporâneo no campo da Educação, pouco foi comentado.
Vale citar que algumas respostas não fizeram qualquer menção ao uso da tecnologia – reconheciam a importância da formação continuada, mas a preocupação recaía tão somente sobre o ensino da norma padrão da língua.
As respostas fracas caracterizavam -se pela generalidade de seu conteúdo ou pelo uso de paráfrases, tendo como base o enunciado da questão, revelando nenhuma ou pouca ideia do estudante sobre os pontos principais.
As respostas boas demonstravam entendimento por parte do estudante sobre o que significa viver num mundo dominado pela tecnologia.
Ficou evidenciado nessas respostas o reconhecimento da necessidade de ajustamento das práticas educacionais a essa nova realidade, por meio de um processo de formação contínua do professor, que lhe permitisse acompanhar e tirar proveito das inovações tecnológicas.
De um modo geral, o estudante de Licenciatura parece estar convencido de que o professor, para alcançar seus objetivos, precisa aperfeiçoar -se continuamente.
Reconhece que, para tanto, há necessidade de que sejam revistas as práticas pedagógicas, ajustando -as a um mundo dominado pelas inovações tecnológicas.
Indicações mais precisas de como tal ajuste deve ser promovido foram, contudo, escassamente observadas.
Também foram registradas, em menor número, respostas que descrevessem com propriedade os reflexos da formação contínua do professor na prática escolar.
Essas lacunas podem estar a revelar a pouca profundidade com que esse tema – tão relevante para a área de Letras – é tratado pelos cursos de graduação.

% matched lemmas: intermédio
\end{textsample}
